% paralleles.tex — f166v, vv. 337–379 : loci paralleli. Généré par .travail/f166v_lp_redac/rediger.py
% à partir des collectes vérifiées (paralleles_337-358.json, paralleles_359-379.json) et des parallèles
% d'Austin–Olson 2004 retrouvés dans les sources ; règles : edition/CONVENTIONS.md § 4
\lp{336–337 ᾿πικηρυκεύεται Εὐριπίδῃ Μήδοις τ᾽]}{Ar.~\textit{Thesm.}~1160–1163 γυναῖκες εἰ βούλεσθε τὸν λοιπόν χρόνον / σπονδὰς ποιήσασθαι πρὸς ἐμέ, νυνὶ πάρα, / ἐφ᾽ ᾧτ᾽ ἀκοῦσαι μηδὲν ὑπ᾽ ἐμοῦ μηδαμὰ / κακὸν τὸ λοιπόν. ταῦτ᾽ ἐπικηρυκεύομαι ; Hdt.~1.60.2 περιελαυνόμενος δὲ τῇ στάσι ὁ Μεγακλέης ἐπεκηρυκεύετο Πεισιστράτῳ, εἰ βούλοιτό οἱ τὴν θυγατέρα ἔχειν γυναῖκα ἐπὶ τῇ τυραννίδι ; Isoc.~4.157 ἐν δὲ τοῖς συλλόγοις ἔτι καὶ νῦν ἀρὰς ποιοῦνται, πρὶν ἄλλο τι χρηματίζειν, εἴ τις ἐπικηρυκεύεται Πέρσαις τῶν πολιτῶν ; Plu.~\textit{Arist.}~10.6 ἔτι δὲ ἀρὰς θέσθαι τοὺς ἱερεῖς ἔγραψεν, εἴ τις ἐπικηρυκεύσαιτο Μήδοις ἢ τὴν συμμαχίαν ἀπολίποι τῶν Ἑλλήνων ; Suid.~ε~2375 Ἐπικηρυκεύεται: διὰ κήρυκός τινος συνθήκας τινὰς ἢ πρεσβείας ποιεῖται. ἢ φιλίαν σπένδει ; Suid.~ε~2376 Ἐπικηρυκεύω· δοτικῇ ; Hdt.~6.97.2 ταῦτα μὲν ἐπεκηρυκεύσατο τοῖσι Δηλίοισι ; Th.~4.27.2 ἔχοντάς τι ἰσχυρὸν αὐτοὺς ἐνόμιζον οὐκέτι σφίσιν ἐπικηρυκεύεσθαι ; Hsch.~ε~4853 ἐπικηρυκεύεται· προκηρύσσεται (Thuc. 7, 83, 2) ; Harp. s.v.~ἐπικηρυκεία ἘΠΙΚΗΡΥΚΕΙΑ τὸ περὶ φιλίας καὶ διαλλαγῶν κήρυκας πέμπειν· Δημοσθένης Φιλιππικῶν ε ; Hdt.~1.69.3 Κροῖσος μὲν δὴ ταῦτα δι᾽ ἀγγέλων ἐπεκηρυκεύετο ; Hdt.~4.80.4 ταῦτά οἱ πέμψας ὁ Σιτάλκης ἐπεκηρυκεύετο ; Ammon.~\textit{Diff.}~184 ἐπικηρύξαι καὶ ἐπικηρυκεύσασθαι διαφέρει … ἐπικηρυκεύσασθαι δὲ ἔλεγον τὸ περὶ συμβιβάσεως καὶ σπονδῶν τὰς πόλεις κήρυκας πρὸς ἀλλήλας διαπέμπεσθαι}
\lp{337–349 ἢ … ἢ … κακῶς ἀπολέσθαι (catalogue de la malédiction)]}{Ar.~\textit{Ran.}~147–151 εἴ που ξένον τις ἠδίκησε πώποτε, / ἢ παῖδα κινῶν τἀργύριον ὑφείλετο, / ἢ μητέρ᾽ ἠλόασεν, ἢ πατρὸς γνάθον / ἐπάταξεν, ἢ ᾽πίορκον ὅρκον ὤμοσεν, / ἢ Μορσίμου τις ῥῆσιν ἐξεγράψατο ; Ar.~\textit{Ran.}~359–362 ἢ στάσιν ἐχθρὰν μὴ καταλύει μηδ᾽ εὔκολός ἐστι πολίταις, / ἀλλ᾽ ἀνεγείρει καὶ ῥιπίζει κερδῶν ἰδίων ἐπιθυμῶν, / ἢ τῆς πόλεως χειμαζομένης ἄρχων καταδωροδοκεῖται, / ἢ προδίδωσιν φρούριον ἢ ναῦς, ἢ τἀπόρρητ᾽ ἀποπέμπει ; SIG³~37–38 A~1–11 ὅστις : φάρμακα : δηλητήρια : ποιοῖ : ἐπὶ Τήιοισιν : τὸ ξυνὸν : ἢ ἐπ᾽ ἰδιώτηι, : κε͂νον : ἀπόλλυσθαι : καὶ αὐτὸν : καὶ γένος : τὸ κένο. … ὅστις : ἐς γῆν : τὴν Τήιην : κωλύοι : σῖτον : ἐσάγεσθαι : … κε͂νον : ἀπόλλυσθαι : καὶ αὐτὸν : καὶ γένος : τὸ κένο ; Ar.~\textit{Thesm.}~1167–1169 ἢν δὲ μὴ πίθησθέ μοι, / ἃ νῦν ὑποικουρεῖτε τοῖσιν ἀνδράσιν / ἀπὸ τῆς στρατιᾶς παροῦσιν ὑμῶν διαβαλῶ ; SIG³~167.12–16 καὶ πρόσθετα ποιήσαντες Μαυσσώλλωι, ἐπαρὰς ἐποιήσαντο περὶ τούτων μήτε προτιθέναι ἔτι παρὰ ταῦτα μηδένα μήτε ἐπιψηφίζειν· εἰ δέ τις ταῦτα παραβαίνοι, ἐξώλη γίνεσθαι καὶ αὐτὸν καὶ τοὺς ἐκείνου πάντας ; Lycurg.~31 τῶν δὲ δικαίως τὰς κρίσεις ἐνισταμένων καὶ τοὺς ἐνόχους ταῖς ἀραῖς ἀκριβῶς ἀποδεικνύντων}
\lp{337 Εὐριπίδῃ Μήδοις τ᾽ ἐπὶ βλάβῃ τινί]}{\textit{Test.} Satyr.~\textit{Vit.\,Eur.} (P.Oxy. IX~1176) fr.~39 col. xii~8–16 εἴ τις ἐπιβουλεύει τι τῷ δήμῳ κακὸν τῷ τῶν γυναικῶν ἢ ᾽πικηρυκεύετ[αι] Εὐριπίδ⟨ῃ Μήδοις⟩ τ᾽ ἐπὶ βλάβῃ τινί ; Ar.~\textit{Lys.}~283 τασδὶ δὲ τὰς Εὐριπίδῃ θεοῖς τε πᾶσιν ἐχθρὰς}
\lp{337 Εὐριπίδῃ]}{Ar.~\textit{Lys.}~368–369 οὐκ ἔστ᾽ ἀνὴρ Εὐριπίδου σοφώτερος ποιητής· οὐδὲν γὰρ οὕτω θρέμμ᾽ ἀναιδές ἐστιν ὡς γυναῖκες ; Ar.~\textit{Thesm.}~85 ὁτιὴ τραγῳδῶ καὶ κακῶς αὐτὰς λέγω ; Vita Eur.~6 (Σ Eur.~\textit{I,} p.~6,~7–11 Schwartz) λέγουσι δὲ καὶ ὅτι αἱ γυναῖκες διὰ τοὺς ψόγους, οὓς ἐποίει εἰς αὐτὰς διὰ τῶν ποιημάτων, τοῖς Θεσμοφορίοις ἐπέστησαν αὐτῷ βουλόμεναι ἀνελεῖν. ἐφείσαντο δὲ αὐτοῦ πρῶτον μὲν διὰ τὰς Μούσας, ἔπειτα δὲ βεβαιωσαμένου μηκέτι αὐτὰς κακῶς ἐρεῖν ; Ar.~\textit{Thesm.}~82–84 αἱ γὰρ γυναῖκες ἐπιβεβουλεύκασί μοι / κἀν Θεσμοφόροιν μέλλουσι περί μου τήμερον / ἐκκλησιάζειν ἐπ᾽ ὀλέθρῳ ; Ar.~\textit{Thesm.}~181–182 μέλλουσί μ᾽ αἱ γυναῖκες ἀπολεῖν τήμερον / τοῖς Θεσμοφορίοις, ὅτι κακῶς αὐτὰς λέγω ; Ar.~\textit{Thesm.}~377–379 καὶ χρηματίζειν πρῶτα περὶ Εὐριπίδου, / ὅ τι χρὴ παθεῖν ἐκεῖνον· ἀδικεῖν γὰρ δοκεῖ / ἡμῖν ἁπάσαις ; Gell.~15.20.6–7 Mulieres fere omnes in maiorem modum exosus fuisse dicitur … Eius odii in mulieres Aristophanes quoque meminit}
\lp{337 Μήδοις]}{Ar.~\textit{Eq.}~478 καὶ πάνθ᾽ ἃ Μήδοις καὶ βασιλεῖ ξυνόμνυτε ; Ar.~\textit{Pax}~107–108 γράψομαι Μήδοισιν αὐτὸν προδιδόναι τὴν Ἑλλάδα ; Ar.~\textit{Thesm.}~365 ἢ Μήδους ἐπάγουσι ; Th.~6.59.4 ὅθεν καὶ ὁρμώμενος ἐς Μαραθῶνα ὕστερον ἔτει εἰκοστῷ ἤδη γέρων ὢν μετὰ Μήδων ἐστράτευσεν ; Σ Ar.~\textit{Thesm.}~365 (Dübner, p.~267,~6–8) ἀντὶ τοῦ εἰπεῖν πολεμίους εἶπε Μήδους ; Suid.~μ~882 Μηδίζω: τὰ τῶν Μήδων φρονῶ. καὶ Μηδίσαντες ; Ar.~\textit{Vesp.}~12 Μῆδός τις ἐπὶ τὰ βλέφαρα νυστακτὴς ὕπνος ; Th.~8.53.1 ὡς ἐξείη αὐτοῖς Ἀλκιβιάδην καταγαγοῦσι καὶ μὴ τὸν αὐτὸν τρόπον δημοκρατουμένοις βασιλέα τε ξύμμαχον ἔχειν ; Suid.~ε~979 ἐμήδισαν: τὰ τῶν Μήδων εἵλοντο ; Suid.~μ~886 Μηδισμὸς ὁ Παυσανίου: ἁλόντος γὰρ ἐπὶ μηδισμῷ, συναιτιῶνται Θεμιστοκλέα ; Zen.~6.16 Τάδε Μῆδος οὐ φυλάξει ; \textit{idem} Erasm.~\textit{Adag.}~2.10.67 (praem. His Medus non insidiabitur ; add. i. e. Haec Medus non observabit) ; App. Prov.~3.92 Μήδων ἀμελῶ: ἐπὶ τῶν καταφρονουμένων}
\lp{337 ἐπὶ βλάβῃ]}{Ar.~\textit{Thesm.}~360 κερδῶν οὕνεκ᾽ ἐπὶ βλάβῃ ; Th.~8.72.1 ὡς οὐκ ἐπὶ βλάβῃ τῆς πόλεως καὶ τῶν πολιτῶν ἡ ὀλιγαρχία κατέστη, ἀλλ᾽ ἐπὶ σωτηρίᾳ τῶν ξυμπάντων πραγμάτων ; D.~21.113 ἐάν τις Ἀθηναίων λαμβάνῃ παρά τινος, ἢ αὐτὸς διδῷ ἑτέρῳ, ἢ διαφθείρῃ τινὰς ἐπαγγελλόμενος, ἐπὶ βλάβῃ τοῦ δήμου ἢ ἰδίᾳ τινὸς τῶν πολιτῶν ; Eur.~\textit{Hipp.}~511 ἅ σ᾽ οὔτ᾽ ἐπ᾽ αἰσχροῖς οὔτ᾽ ἐπὶ βλάβῃ φρενῶν ; Th.~1.40.2 οὐ τοῖς ἐπὶ βλάβῃ ἑτέρων ἰοῦσιν ἡ ξυνθήκη ἐστίν ; Pl.~\textit{R.}~334b ἐπ᾽ ὠφελίᾳ μέντοι τῶν φίλων καὶ ἐπὶ βλάβῃ τῶν ἐχθρῶν ; X.~\textit{Mem.}~4.2.19 τῶν δὲ δὴ τοὺς φίλους ἐξαπατώντων ἐπὶ βλάβῃ ; Isoc.~4.130 χρὴ δὲ κατηγορεῖν μὲν ἡγεῖσθαι τοὺς ἐπὶ βλάβῃ τοιαῦτα λέγοντας, νουθετεῖν δὲ τοὺς ἐπ᾽ ὠφελείᾳ λοιδοροῦντας ; Isoc.~8.72 τοὺς μὲν ἐπὶ βλάβῃ λοιδοροῦντας μισεῖν ὡς κακόνους ὄντας τῇ πόλει ; D.~24.204 ὃς νόμον εἰσενήνοχ᾽ ἐπὶ βλάβῃ τοῦ πλήθους ; D.~37.26 κἄπειτα πείσας τοὺς οἰκέτας τοὺς ἐμοὺς καθέζεσθαι εἰς τὸν κεγχρεῶνα ἐπὶ βλάβῃ τῇ ἐμῇ ; Hsch.~β~649 βλάβη· λύμη. φθορά (Eur. Hipp. 511) ; Zen.~3.28 Δίκη δίκην ἔτικτε καὶ βλάβην βλάβη: ἐπὶ τῶν φιλοδίκων καὶ συνειρόντων δίκας δίκαις ; \textit{idem} Greg. Cypr.~2.14 (βλάβη βλάβην ; δίκαις δίκας ; add. καὶ ἐπὶ τῶν πυκναῖς προσπιπτόντων συμφοραῖς) ; Pl.~\textit{Euthphr.}~13c ἢ ἐπὶ βλάβῃ οἴει τοῦ θεραπευομένου τὴν θεραπείαν εἶναι ; Pl.~\textit{Lg.}~933d ὃς ἂν φαρμακεύῃ τινὰ ἐπὶ βλάβῃ μὴ θανασίμῳ ; X.~\textit{Mem.}~2.3.19 τοῖς ἐπ᾽ ὠφελείᾳ πεποιημένοις ἐπὶ βλάβῃ χρῆσθαι ; Isoc.~12.219 ὅσοι τοῖς πράγμασι τοῖς εὑρημένοις ἐπ᾽ ὠφελίᾳ, τούτοις ἐπὶ βλάβῃ χρώμενοι τυγχάνουσι ; Aeschin.~3.228 καὶ τὴν φύσιν μου γεγενῆσθαι ἐπὶ βλάβῃ τῶν ἀκουόντων ; Arist.~\textit{Rh.}~1374a οὐδὲ πάντως, εἰ λάθρᾳ ἔλαβεν, ἔκλεψεν, ἀλλ᾽ εἰ ἐπὶ βλάβῃ ; Hsch.~λ~1414 λύμη· βλάβη, φθορά ; Harp. s.v.~παραφρυκτωρεῖν ἐπὶ συμφέροντι μὲν τῶν πολεμίων, ἐπὶ βλάβῃ δὲ τῶν φιλίων}
\lp{337 ἐπὶ βλάβῃ (prosodie : ἐπῐ̄ βλ-)]}{Ar.~\textit{Ran.}~1049 καὶ τί βλάπτουσ᾽ ὦ σχέτλι᾽ ἀνδρῶν τὴν πόλιν ; Ar.~\textit{Ran.}~1151 σὺ δ᾽ ἐπιτήρει τὸ βλάβος ; Ar.~\textit{Ran.}~1428 βραδὺς πέφυκε μεγάλα δὲ βλάπτειν ταχύς}
\lp{337 ἐπὶ βλάβῃ τινί]}{X.~\textit{Cyr.}~2.2.12 μήτε ἐπὶ τῷ αὑτῶν κέρδει μήτ᾽ ἐπὶ ζημίᾳ τῶν ἀκουόντων μήτε ἐπὶ βλάβῃ μηδεμιᾷ ; Suid.~κ~1738 λέγεται δὲ κλέπτης καὶ ὁ ἐπὶ βλάβῃ τινὶ βλάψας ; Suid.~ο~411 Ὀνυχίζεται: ἀκριβολογεῖται. οὕτως Ἀριστοφάνης. ἐν δὲ τῇ συνηθείᾳ τὸ ἐπὶ βλάβῃ τινι ἐξαπατηθῆναι ; Eur.~\textit{HF}~881–882 ἅρμασι δ᾽ ἐνδίδωσι / κέντρον ὡς ἐπὶ λώβᾳ ; Soph.~\textit{Aj.}~143 μεγάλοι θόρυβοι κατέχουσ᾽ ἡμᾶς / ἐπὶ δυσκλείᾳ ; Soph.~\textit{Ant.}~791–792 σὺ καὶ δικαίων ἀδίκους φρένας παρασπᾷς ἐπὶ λώβᾳ ; Hdt.~3.14.1 κατίσας ἐς τὸ προάστειον ἐπὶ λύμῃ τὸν βασιλέα τῶν Αἰγυπτίων Ψαμμήνιτον}
\lp{337–338 ἐπὶ βλάβῃ τινὶ τῇ τῶν γυναικῶν]}{Pl.~\textit{R.}~590d οὐκ ἐπὶ βλάβῃ τῇ τοῦ δούλου οἰόμενοι δεῖν ἄρχεσθαι αὐτόν}
\lp{338–339 ἢ τυραννεῖν ἐπινοεῖ ἢ τὸν τύραννον συγκατάγειν]}{And.~1.97 καὶ ἐάν τις τυραννεῖν ἐπαναστῇ ἢ τὸν τύραννον συγκαταστήσῃ ; D.~20.107 παρὰ δ᾽ ἡμῖν ταύτης μὲν ὁ δῆμος κύριος, καὶ ἀραὶ καὶ νόμοι καὶ φυλακαὶ ὅπως μηδεὶς ἄλλος κύριος γενήσεται ; Arist.~\textit{Ath.}~16.10 θέσμια τάδε Ἀθηναίων καὶ πάτρια· ἐάν τινες τυραννεῖν ἐπανιστῶνται [ἐπὶ τυραννίδι] ἢ συγκαθιστῇ τὴν τυραννίδα, ἄτιμον εἶναι καὶ αὐτὸν καὶ γένος ; Lycurg.~125 ἐψηφίσαντο γὰρ καὶ ὤμοσαν, ἐάν τις τυραννίδι ἐπιτιθῆται ἢ τὴν πόλιν προδιδῷ ἢ τὸν δῆμον καταλύῃ, τὸν αἰσθανόμενον καθαρὸν εἶναι ἀποκτείναντα ; Σ Ar.~\textit{Thesm.}~339 (Dübner, p.~266,~40–42) ἐκ τῆς κατάρας τῆς ἐπὶ τοῖς Πεισιστρατίδαις γενομένης ταῦτα παρέμιξεν. ἔοικέ τινα παραμιγνύναι τῶν ἐφ᾽ Ἱππίᾳ ψηφισθέντων ; Poll.~6.151–152 (Bethe II, p.~40,~10–21) Ἐξ ἀδικημάτων ὀνόματα, ἐφ᾽ οἷς εἰσι δίκαι καὶ γραφαί. κλέπτης … μοιχός, ἀνδροφόνος, … προδότης, … τύραννος, ἀνδραποδιστής … ἃ δ᾽ οὐκ ἔστιν ἑνὶ ὀνόματι εἰπεῖν ἀλλὰ μετοχαῖς, βουλεύων τυραννίδα, καταλύων τὴν πολιτείαν, ἐπὶ καταλύσει τοῦ δήμου συνιών, τὰ δικαστήρια δεκάζων, ἐξαπατῶν τὸν δῆμον ἢ τὴν βουλὴν ἢ τὸ δικαστήριον, τὰ ψευδῆ μαρτυρῶν ; Zen.~4.65 Κανθάρου σοφώτερος: ἡ παροιμία εἴρηται ἐπὶ τῶν πονηρῶν καὶ κακούργων· ἐπειδή τις Ἀθήνησι κάπηλος Κάνθαρος καλούμενος, ἐπὶ πονηρίᾳ καὶ προδοσίᾳ θανάτου ἐζημιώθη· [οἱ δὲ κάπηλον ὄντα φασὶν ἀπερισκέπτως ἐπιθέσθαι τυραννίδι καὶ ἀπολέσθαι ; Th.~6.60.1 καὶ πάντα αὐτοῖς ἐδόκει ἐπὶ ξυνωμοσίᾳ ὀλιγαρχικῇ καὶ τυραννικῇ πεπρᾶχθαι ; D.~24.149 καὶ τύραννον οὐ ψηφιοῦμαι εἶναι οὐδ᾽ ὀλιγαρχίαν ; Arist.~\textit{Ath.}~8.4 καὶ τοὺς ἐπὶ καταλύσει τοῦ δήμου συνισταμένους ἔκρινεν, Σόλωνος θέντος νόμον εἰσαγγελίας περὶ αὐτῶν ; Hyp.~4.7–8 ἐάν τις, φησί, τὸν δῆμον τὸν Ἀθηναίων καταλύῃ ; Th.~6.28.2 ἐμεγάλυνον καὶ ἐβόων ὡς ἐπὶ δήμου καταλύσει τά τε μυστικὰ καὶ ἡ τῶν Ἑρμῶν περικοπὴ γένοιτο}
\lp{338 ἢ τυραννεῖν ἐπινοεῖ]}{Ar.~\textit{Vesp.}~463–464 ἆρα δῆτ᾽ οὐκ αὐτὰ δῆλα τοῖς πένησιν, ἡ τυραννὶς ; Ar.~\textit{Vesp.}~488 ὡς ἅπανθ᾽ ὑμῖν τυραννίς ἐστι καὶ ξυνωμόται ; Ar.~\textit{Lys.}~630 ἀλλὰ ταῦθ᾽ ὕφηναν ἡμῖν ἄνδρες ἐπὶ τυραννίδι ; Ar.~\textit{Lys.}~631–632 ἀλλ᾽ ἐμοῦ μὲν οὐ τυραννεύσουσ᾽, ἐπεὶ φυλάξομαι / καὶ φορήσω τὸ ξίφος τὸ λοιπὸν ἐν μύρτου κλαδί ; Ar.~\textit{Thesm.}~1143–1144 φάνηθ᾽ ὦ τυράννους στυγοῦσ᾽ ὥσπερ εἰκός ; Eur.~\textit{Phoen.}~524–525 εἴπερ γὰρ ἀδικεῖν χρή, τυραννίδος πέρι κάλλιστον ἀδικεῖν, τἄλλα δ᾽ εὐσεβεῖν χρεών ; Eur.~\textit{Phoen.}~560–561 πότερα τυραννεῖν ἢ πόλιν σῷσαι θέλεις, ἐρεῖς τυραννεῖν ; Eur.~\textit{Supp.}~429 οὐδὲν τυράννου δυσμενέστερον πόλει ; Eur.~\textit{Tro.}~1169 γάμων τε καὶ τῆς ἰσοθέου τυραννίδος ; Aesch.~\textit{Ag.}~1354–1355 φροιμιάζονται γὰρ ὡς / τυραννίδος σημεῖα πράσσοντες πόλει ; Soph.~\textit{OT}~541–542 ἄνευ τε πλήθους καὶ φίλων τυραννίδα / θηρᾶν ; Th.~6.53.3 ἐπιστάμενος γὰρ ὁ δῆμος ἀκοῇ τὴν Πεισιστράτου καὶ τῶν παίδων τυραννίδα χαλεπὴν τελευτῶσαν γενομένην … ἐφοβεῖτο αἰεὶ καὶ πάντα ὑπόπτως ἐλάμβανεν ; Pl.~\textit{R.}~568a–b ὁ Εὐριπίδης διαφέρων ἐν αὐτῇ … ὡς ἄρα σοφοὶ τύραννοί εἰσι τῶν σοφῶν συνουσίᾳ … καὶ ὡς ἰσόθεόν γ᾽, ἔφη, τὴν τυραννίδα ἐγκωμιάζει ; \textit{idem} Zen.~5.98 (om. εἰσι ; add. τοῦτο Σοφοκλέους ἐστὶν ἐξ Αἴαντος τοῦ Λοκροῦ· Πλάτων δὲ Εὐριπίδου φησὶν εἶναι τὸ ἰαμβεῖον), App. Prov.~4.76 (om. εἰσι), Suid.~σ~823 (om. εἰσι ; τῇ τῶν σοφῶν) ; App. Prov.~2.99 Ἡ γὰρ τυραννὶς ἀδικίας μήτηρ ἔφυ ; Ar.~\textit{Vesp.}~417 ταῦτα δῆτ᾽ οὐ δεινὰ καὶ τυραννίς ἐστιν ἐμφανής ; Ar.~\textit{Vesp.}~473 ὦ μισόδημε καὶ μοναρχίας ἐραστά ; Ar.~\textit{Vesp.}~487 ὅστις ἡμῶν ἐπὶ τυραννίδ᾽ ὧδ᾽ ἐστάλης ; Ar.~\textit{Vesp.}~495 οὗτος ὀψωνεῖν ἔοιχ᾽ ἅνθρωπος ἐπὶ τυραννίδι ; Eur.~\textit{Ion}~621–622 τυραννίδος δὲ τῆς μάτην αἰνουμένης τὸ μὲν πρόσωπον ἡδύ, τἀν δόμοισι δὲ ; Soph.~\textit{OT}~873 ὕβρις φυτεύει τύραννον ; Sol.~fr.~9.3–4~W. ἀνδρῶν δ᾽ ἐκ μεγάλων πόλις ὄλλυται· ἐς δὲ μονάρχου / δῆμος ἀιδρίῃ δουλοσύνην ἔπεσεν ; Hdt.~1.59.3 καταφρονήσας τὴν τυραννίδα ἤγειρε τρίτην στάσιν ; Hdt.~3.53.4 τυραννὶς χρῆμα σφαλερόν, πολλοὶ δὲ αὐτῆς ἐρασταί εἰσι ; Hdt.~5.71.1 ἦν Κύλων τῶν Ἀθηναίων ἀνὴρ Ὀλυμπιονίκης· οὗτος ἐπὶ τυραννίδι ἐκόμησε ; Th.~1.126.5 κατέλαβε τὴν ἀκρόπολιν ὡς ἐπὶ τυραννίδι ; Th.~6.15.4 φοβηθέντες γὰρ αὐτοῦ οἱ πολλοὶ τὸ μέγεθος τῆς τε κατὰ τὸ ἑαυτοῦ σῶμα παρανομίας ἐς τὴν δίαιταν καὶ τῆς διανοίας … ὡς τυραννίδος ἐπιθυμοῦντι πολέμιοι καθέστασαν ; Plu.~\textit{Sol.}~19.3 ἐπὶ φόνῳ ἢ σφαγαῖσιν ἢ ἐπὶ τυραννίδι ἔφευγον ; Hsch.~μ~1604 μονάρχου· τυράννου (Sol. fr. 9, 3 W.) ; Suid.~π~80 Παλλήνη δὲ δῆμος τῆς Ἀττικῆς, ἔνθα Πεισιστράτῳ βουλομένῳ τυραννεῖν καὶ Ἀθηναίοις μὴ βουλομένοις συνέστη πόλεμος ; Diogenian.~2.71 Ἀρχαιότερος Ἰβύκου: ἐπὶ τῶν εὐηθῶν. Οὗτος γὰρ τυραννεῖν δυνάμενος ἀπεδήμησεν ; Diogenian.~7.77b Πόῤῥω Διός τε καὶ κεραυνοῦ: παραινετικὴ ὅτι δεῖ φεύγειν τοὺς τυράννους ὡς ἀπὸ κεραυνοῦ ; Tosi n°~134 (Ἀνδρὸς ⟨δ᾽⟩ ὑπ᾽ ἐσθλοῦ καὶ τυραννεῖσθαι καλόν) Ἀνδρὸς ⟨δ᾽⟩ ὑπ᾽ ἐσθλοῦ καὶ τυραννεῖσθαι καλόν ; Ar.~\textit{Eq.}~1111–1114 ὦ Δῆμε καλήν γ᾽ ἔχεις / ἀρχήν, ὅτε πάντες ἄνθρωποι / δεδίασί σ᾽ ὥσπερ / ἄνδρα τύραννον ; Eur.~\textit{Hipp.}~1013 ἀλλ᾽ ὡς τυραννεῖν ἡδύ ; Hsch.~α~4235 ἀναγκομόναρχος· ὁ τύραννος ; Hsch.~α~5329 ἀντάρτης· τύραννος, ἀντιβαίνων βασιλεῖ ; Hsch.~λ~134 λαίλας· ὁ ⟨μὴ ἐκ γένους⟩ τύραννος, ὑπὸ Λυδῶν ; Hsch.~τ~1666 τυραννοδαίμονα· ἣν οὐκ ἄν τις τύραννον μόνον εἴποι, ἀλλὰ καὶ δαίμονα (com. ad. fr. 433 K.–A.) ; Hsch.~τ~1667 τύραννος· παράνομος βασιλεύς, ἢ ἄρχων ἀπηνὴς καὶ ἀπάνθρωπος ; Suid.~β~144 διαφέρει δὲ βασιλεὺς τυράννου}
\lp{338 ἐπινοεῖ]}{Ar.~\textit{Eq.}~1202 εἴπ᾽, ἀντιβολῶ, πῶς ἐπενόησας ἁρπάσαι ; Ar.~\textit{Nub.}~1039–1040 ὅτι πρώτιστος ἐπενόησα / τοῖσιν νόμοις καὶ ταῖς δίκαις τἀναντί᾽ ἀντιλέξαι ; Ar.~\textit{Pax}~1268 ἀλλ᾽ ὅ τι περ ᾄδειν ἐπινοεῖς ὦ παιδίον ; Ar.~\textit{Thesm.}~766 τίς πεῖρα, τίς ἐπίνοι᾽ ; Hermipp.~fr.~4~K.-A. (2~Kock) ἀπὸ τῆς τραπέζης τουτονὶ τὸν στήμονα / ἄττεσθ᾽ ἐπινοῶν ; [Eur.]~\textit{Rh.}~195 μέγας ἀγών, μεγάλα δ᾽ ἐπινοεῖς ἑλεῖν ; Hdt.~1.27.1 τὸ ἐνθεῦτεν ἐπενόεε νέας ποιησάμενος ἐπιχειρέειν τοῖσι νησιώτῃσι ; Hdt.~3.31.2 ὅτι οὐκ ἐωθότα ἐπενόεε ποιήσειν ; Th.~2.8.1 ὀλίγον τε ἐπενόουν οὐδὲν ἀμφότεροι, ἀλλ᾽ ἔρρωντο ἐς τὸν πόλεμον ; Pl.~\textit{Cri.}~52a εἴπερ ποιήσεις ἃ ἐπινοεῖς ; Lys.~6.31 ἃ τούτῳ ὁ θεὸς … ἐπινοεῖν δίδωσιν, ἀλλὰ τιμωρούμενος τῶν γεγενημένων ἀσεβημάτων}
\lp{338 τῇ τῶν γυναικῶν]}{Ar.~\textit{Thesm.}~306–307 τὸν δῆμον τὸν Ἀθηναίων καὶ τὸν τῶν γυναικῶν ; Ar.~\textit{Thesm.}~373–374 ἔδοξε τῇ βουλῇ τάδε τῇ τῶν γυναικῶν ; Ar.~\textit{Thesm.}~1145 δῆμός τοί σε καλεῖ γυναικῶν}
\lp{339 ἢ τὸν τύραννον συγκατάγειν]}{Ar.~\textit{Vesp.}~502 ἤρετ᾽ εἰ τὴν Ἱππίου καθίσταμαι τυραννίδα ; Ar.~\textit{Av.}~1074–1075 ἤν τε τῶν τυράννων τίς τινα τῶν τεθνηκότων ἀποκτείνῃ, τάλαντον λαμβάνειν ; Ar.~\textit{Lys.}~618 καὶ μάλιστ᾽ ὀσφραίνομαι τῆς Ἱππίου τυραννίδος ; Hdt.~1.60.5 ἔς τε τοὺς δήμους φάτις ἀπίκετο ὡς Ἀθηναίη Πεισίστρατον κατάγει, καὶ οἱ ἐν τῷ ἄστεϊ πειθόμενοι τὴν γυναῖκα εἶναι αὐτὴν τὴν θεὸν προσεύχοντό τε τὴν ἄνθρωπον καὶ ἐδέκοντο Πεισίστρατον ; Hdt.~5.91.1 μετεπέμποντο Ἱππίην τὸν Πεισιστράτου ἀπὸ Σιγείου ; Hdt.~5.96.2 ὁ δὲ Ἀρταφρένης ἐκέλευε σφέας, εἰ βουλοίατο σόοι εἶναι, καταδέκεσθαι ὀπίσω Ἱππίην ; Hdt.~6.107.1 τοῖσι δὲ βαρβάροισι κατηγέετο Ἱππίης ὁ Πεισιστράτου ἐς τὸν Μαραθῶνα ; Hdt.~6.121.1 βουλομένους ὑπὸ βαρβάροισί τε εἶναι Ἀθηναίους καὶ ὑπὸ Ἱππίῃ ; Isoc.~16.13 ἐπὶ τίνα δ᾽ ἢ πόλιν ἢ φίλον ἢ ξένον οὐκ ἤλθετε δεησόμενοι συγκαταγαγεῖν ὑμᾶς ; D.~17.3 εἰ ἦσαν κατὰ τὸν νυνὶ χρόνον οἱ Πεισιστρατίδαι καί τις ἐβιάζετο κατάγειν αὐτοὺς δευρί, ἁρπάσαντας ἂν ὑμᾶς τὰ ὅπλα πάντα κίνδυνον ὑπομεῖναι ; Arist.~\textit{Ath.}~14.4 πάλιν ἐπικηρυκευσάμενος πρὸς τὸν Πεισίστρατον, ἐφ᾽ ᾧ τε τὴν θυγατέρα αὐτοῦ λήψεται, κατήγαγεν αὐτὸν … προδιασπείρας γὰρ λόγον, ὡς τῆς Ἀθηνᾶς καταγούσης Πεισίστρατον ; Plu.~\textit{De\,Herod.\,malign.}~858c τοὺς Ἀλκμεωνίδας … εἰς αἰτίαν ἐμβαλὼν προδοσίας δέξασθαί φησι τὸν Πεισίστρατον ἐκ τῆς φυγῆς καὶ συγκαταγαγεῖν ἐπὶ τῷ γάμῳ τῆς Μεγακλέους θυγατρός ; Hdt.~5.92η.5 ὅτε ὑμέας εἴδομεν μεταπεμπομένους Ἱππίην … μὴ κατιστάναι τυραννίδας ἐς τὰς πόλις ; D.~12.7 τοῖς Πεισιστρατίδαις ἐπετίμων ὡς ἐπάγουσι τὸν Πέρσην ἐπὶ τοὺς Ἕλληνας ; Plu.~\textit{De\,Herod.\,malign.}~860d εἶτα μεταπεμψαμένους Ἱππίαν ἀπὸ Σιγείου κατάγειν εἰς τὰς Ἀθήνας ; Paus.~2.21.8 φυγόντα δὲ ἐς Σπάρτην Λακεδαιμόνιοι κατάγειν ἐπειρῶντο ἐπὶ τυραννίδι ; Ath.~13.609c καὶ τὴν καταγαγοῦσαν δὲ Πεισίστρατον ἐπὶ τὴν τυραννίδα, ὡς Ἀθηνᾶς Παλληνίδος εἶδος ἔχουσαν, καλήν φησι γεγονέναι ; App. Prov.~2.64 Ἐν μύρτου κλαδὶ τὸ ξίφος κρατήσω: … Ἁρμόδιος γὰρ καὶ Ἀριστογείτων κτεῖναι θέλοντες Ἵππαρχον τὸν τύραννον, ἐν μύρτοις τὸ ξίφος ῥίψαντες, τὸν φόνον εἰργάσαντο ; Hdt.~7.6.2 τοῦτο δὲ Πεισιστρατιδέων οἱ ἀναβεβηκότες ἐς Σοῦσα, τῶν τε αὐτῶν λόγων ἐχόμενοι τῶν καὶ οἱ Ἀλευάδαι ; Hdt.~8.52.2 οὐδὲ λόγους τῶν Πεισιστρατιδέων προσφερόντων περὶ ὁμολογίης ἐνεδέκοντο ; Hsch.~β~1322 Βυρσίνης· Μυρσίνης· ἦν δὲ αὕτη γυνὴ Ἱππίου ⟨τοῦ⟩ τυράννου. ἐν τοῖς Ἱππεῦσι (Ar. Equ. 449) ; Hsch.~κ~1887 κατωνάκη· … δοκοῦσι δὲ τούτο ἀμφιέσασθαι Ἀθηναῖοι τῶν περὶ Πεισίστρατον τυράννων ἐπαναγκασάντων, ἵνα ὑπὸ εὐτελείας μὴ κατίωσιν εἰς τὸ ἄστυ οἱ πολῖται (Ar. Eccl. 724) ; Zen.~4.76 Καὶ σφάκελοι ποιοῦσιν ἀτέλειαν: Πεισίστρατος, ὡς φασὶν, ὁ τύραννος δεκάτην τῶν γεωργουμένων ἀπῄτει τοὺς Ἀθηναίους}
\lp{339 συγκατάγειν]}{Aeschin.~2.147 συγκατάγειν δὲ τὸν δῆμον ; Aeschin.~2.78 ἔφυγε μὲν ἐπὶ τῶν τριάκοντα, συγκατήγαγε δὲ τὸν δῆμον ; D.H.~5.64.3 καὶ συγκαταγάγῃ Ταρκύνιον ἐπὶ τὴν ἀρχήν ; Suid.~κ~495 Κατάγειν καὶ τὸ τοὺς πεφευγότας εἰς τὰς ἑαυτῶν πατρίδας ἀποκαθιστᾶν. ὁ δὲ ἐπειρᾶτο κατάγειν τοὺς ἐκπεπτωκότας ; Pl.~\textit{Ep.}~7.333e καὶ δὴ καὶ τούτω τὼ συγκαταγαγόντε αὐτὸν φίλω}
\lp{339–340 παιδίον ὑποβαλλομένης]}{Ar.~\textit{Pax}~676–678 οὐκ ἦν ἄρ᾽ οὗπέρ φησιν εἶναι τοῦ πατρός. εἰ γάρ ποτ᾽ ἐξέλθοι στρατιώτης, εὐθέως ἀποβολιμαῖος τῶν ὅπλων ἐγίγνετο ; Ar.~\textit{Thesm.}~407–408 εἶεν, γυνή τις ὑποβαλέσθαι βούλεται ἀποροῦσα παίδων ; Ar.~\textit{Thesm.}~502–503 ἑτέραν δ᾽ ἐγᾦδ᾽ ἣ ᾽φασκεν ὠδίνειν γυνὴ δέχ᾽ ἡμέρας, ἕως ἐπρίατο παιδίον ; Ar.~\textit{Thesm.}~564–565 οὐδ᾽ ὡς σὺ τῆς δούλης τεκούσης ἄρρεν εἶτα σαυτῇ τοῦθ᾽ ὑπεβάλου, τὸ σὸν δὲ θυγάτριον παρῆκας αὐτῇ ; Epinic.~\textit{Ὑποβαλλόμεναι} ap. Ath.~11.469a μνημονεύει τοῦ ποτηρίου τούτου καὶ Ἐπίνικος ἐν Ὑποβαλλομέναις ; Plaut.~\textit{Capt.}~1030–1031 neque in hac subigitationes sunt neque ulla amatio / nec pueri suppositio nec argenti circumductio ; Ter.~\textit{Eun.}~38–40 parasitum edacem, gloriosum militem, / puerum supponi, falli per servum senem, / amare, odisse, suspicari? ; Eur.~\textit{Alc.}~636–639 οὐκ ἦσθ᾽ ἄρ᾽ ὀρθῶς τοῦδε σώματος πατήρ; οὐδ᾽ ἡ τεκεῖν φάσκουσα καὶ κεκλημένη μήτηρ μ᾽ ἔτικτε; δουλίου δ᾽ ἀφ᾽ αἵματος μαστῷ γυναικὸς σῆς ὑπεβλήθην λάθρᾳ ; Juv.~6.602–603 Transeo suppositos et gaudia votaque saepe / ad spurcos decepta lacus ; Hdt.~1.137.2 ἤτοι ὑποβολιμαῖα ἐόντα ἢ μοιχίδια ; Hdt.~5.41.2 φάμενοι αὐτὴν κομπέειν ἄλλως βουλομένην ὑποβαλέσθαι. δεινὰ δὲ ποιεύντων αὐτῶν, τοῦ χρόνου συντάμνοντος, ὑπ᾽ ἀπιστίης οἱ ἔφοροι τίκτουσαν τὴν γυναῖκα περιιζόμενοι ἐφύλαξαν ; Pl.~\textit{Menex.}~237e ᾧ καὶ γυνὴ δήλη τεκοῦσά τε ἀληθῶς καὶ μή, ἀλλ᾽ ὑποβαλλομένη, ἐὰν μὴ ἔχῃ πηγὰς τροφῆς τῷ γεννωμένῳ ; D.~21.149 ἡ μὲν γὰρ ὡς ἀληθῶς μήτηρ, ἡ τεκοῦσ᾽ αὐτόν, πλεῖστον ἁπάντων ἀνθρώπων εἶχε νοῦν, ἡ δὲ δοκοῦσα καὶ ὑποβαλομένη πασῶν ἦν ἀνοητοτάτη γυναικῶν ; Hsch.~δ~2716 δωλέγγονος· ὑπόβλητος ; Hsch.~κ~2549 κηρίῳ βύσασα· εἰς τὸ στόμα τῶν παίδων ἐτίθεσαν κηρίον, ὑπὲρ τοῦ μὴ βοᾶν ὁπότε ὑποβάλλοιεν ; Hsch.~υ~576 ὑποβολιμαῖον· οὐ γνήσιον, ἀλλὰ νόθον, ὑποβαλλόμενον, ὡς ἀπὸ τῶν χαμαιριφῶν παιδίων, ἅπερ ἑαυταῖς ὑποβάλλουσιν αἱ γυναῖκες ; \textit{idem} Synag.~υ~117 (Cunningham, p.~485) (ὑποβολιμαῖοι· ὑποβαλλόμενοι καὶ εἰσποιητοί ; ὡς ἐπὶ ; ἑαυτῶν ἀποβάλλουσιν) ; Apostol.~17.58 Ὑποβολιμαῖος εἶ: ἐπὶ τῶν εὐτελῶν· ἀπὸ τῶν ὑποβαλλομένων παιδίων καὶ χαμαιρριφῶν, ἅπερ ἑαυτῶν ἀποβάλλουσιν ἢ ὑποβάλλουσιν αἱ γυναῖκες ; Telecl.~fr.~41~Kock (Plu.~\textit{Nic.}~4.4) Χαρικλέης μὲν οὖν ἔδωκε μνᾶν, ἵν᾽ αὐτὸν μὴ λέγῃ / ὡς ἔφυ τῇ μητρὶ παίδων πρῶτος ἐκ βαλλαντίου ; Alex.~\textit{Ὑποβολιμαῖος} ap. Ath.~11.502b χυτρίδες. Ἄλεξις ἐν Ὑποβολιμαίῳ ; Men.~\textit{Ὑποβολιμαῖος} ap. Suid.~π~959 Μένανδρος Ὑποβολιμαίῳ ἢ Ἀγροίκῳ ; Ter.~\textit{An.}~512–515 nunc, postquam videt / nuptias domi apparari, missa est ancilla ilico / obstetricem arcessitum ad eam, et puerum ut afferret simul ; Soph.~\textit{OT}~779–780 ἀνὴρ γὰρ ἐν δείπνοις μ᾽ ὑπερπλησθεὶς μέθῃ / καλεῖ παρ᾽ οἴνῳ, πλαστὸς ὡς εἴην πατρί ; Pl.~\textit{R.}~537e εἴ τις ὑποβολιμαῖος τραφείη ἐν πολλοῖς μὲν χρήμασι ; D.~9.31 εἰ δέ γε δοῦλος ἢ ὑποβολιμαῖος τὰ μὴ προσήκοντ᾽ ἀπώλλυε καὶ ἐλυμαίνετο ; Plu.~\textit{Mul.\,virt.}~258d συνέπεισεν ἐξ ἑτέρας γυναικὸς παιδοποιησάμενον αὐτῇ τὸ παιδίον περιιδεῖν ὑποβαλλόμενον ; Apollod. Bibl.~3.5.7 ἡ δὲ ἀνελοῦσα ὑποβάλλεται … διὰ φθόνον ὠνειδίζετο ὑπόβλητος ; Paus.~4.9.5 τὴν δὲ γυναῖκα ἣ Λυκίσκῳ συνῴκησεν, ὡς τεκεῖν οὔκουν οἵα τε ἦν, [ἐν τούτῳ] τὴν παῖδα ὑποβαλέσθαι ; Lex. Rhet. (Bekker, Anecd. I, p.~311,~33 –~312,~2) Ὑποβολῆς γραφὴ τί ἐστιν: εἶδος ἐγκλήματος. εἴ τις ἐγκαλοίη τινὶ ὡς ὑποβολιμαῖος εἴη, ἐγράφετο ὑποβολῆς, καὶ ἁλόντα αὐτὸν ἔδει πεπρᾶσθαι ; Suid.~α~2673 ὑποβλήτους ἄν εἶπέ τις αὐτοὺς εἶναι καὶ νόθους ἄντικρυς ; Suid.~υ~455 Ὑπόβλητον: οὐκ ἀληθές, ἀλλὰ πεπλασμένον. ἀπὸ τῶν ὑποβαλλομένων εἰς τὰ γένη τοὺς παῖδας ; Erasm.~\textit{Adag.}~4.3.46 Subdititius es. Ὑποβολιμαῖος εἶ. i. e. Suppositius es. Subdititii sive submissitii vocantur partus adsimulati : inde transfertur in hominem fictum, \& per dolum ab alio quopiam allegatum subornatumque ; X.~\textit{Cyn.}~7.3 ὑπὸ τῇ τεκούσῃ ἐᾶν καὶ μὴ ὑποβάλλειν ὑφ᾽ ἑτέραν κύνα ; Liv.~3.44.9 furtoque inde in domum Vergini translatam suppositam ei esse ; Hsch.~ν~617 νοθογέννητα· πορνειογέννητα. ⟨νόθοι·⟩ οἱ μὴ γνήσιοι παῖδες, ἀλλ᾽ ἀπὸ πόρνης ἢ φίλης ἢ δούλης ἢ παλλακῆς}
\lp{340 κατεῖπε]}{Ar.~\textit{Pax}~376–377 μὴ πρὸς τῶν θεῶν ἡμῶν κατείπῃς, ἀντιβολῶ σε δέσποτα ; Ar.~\textit{Eccl.}~495 μὴ καί τις ὄψεθ᾽ ἡμᾶς χἠμῶν ἴσως κατείπῃ ; Suid.~κ~1034 Κατεῖπε· γενικῇ. κατηγόρησεν. οὐ δικαίως ὀργίζῃ τῇ γυναικί· οὐ γάρ σου κατεῖπεν, ἀλλὰ τοῦ δεσμοφύλακος ; Pl.~\textit{Tht.}~149a μὴ μέντοι μου κατείπῃς πρὸς τοὺς ἄλλους ; Hsch.~κ~1594 κατειπών· διαβαλών ; Phot.~κ~420 Theodoridis κατείρηκεν· κατεῖπεν λάθρα ; Eur.~\textit{Hec.}~243 ἔγνω δέ σ᾽ Ἑλένη καὶ μόνῃ κατεῖπ᾽ ἐμοί ; Hdt.~2.89.2 λαμφθῆναι γὰρ τινὰ φασὶ μισγόμενον νεκρῷ προσφάτῳ γυναικός, κατειπεῖν δὲ τὸν ὁμότεχνον ; Hsch.~κ~1484 καταφάναι· κατειπεῖν ; \textit{idem} Synag.~κ~175 (Cunningham, p.~291)}
\lp{340 δούλη]}{Hsch.~α~176 ἄβρα· δούλη, παλλακή (Men. fr. 411, 3 K.-A.)}
\lp{340–341 δούλη τινὸς … ἐνετρύλλισεν τῷ δεσπότῃ]}{Lys.~1.20 κατηγόρει πρῶτον μὲν ὡς μετὰ τὴν ἐκφορὰν αὐτῇ προσίοι … καὶ ὡς Θεσμοφορίοις ἐμοῦ ἐν ἀγρῷ ὄντος ᾤχετο εἰς τὸ ἱερὸν μετὰ τῆς μητρὸς τῆς ἐκείνου· καὶ τἆλλα τὰ γενόμενα πάντα ἀκριβῶς διηγήσατο ; Hermog.~\textit{Stat.}~3 (p.~45,~5–11 Rabe) εἶτα βασανίζων ὁ πατὴρ τὰς θεραπαινίδας περὶ μὲν τῆς φαρμακείας ἀκήκοε πλέον οὐδέν, μοιχείαν δὲ τῆς γυναικὸς πρὸς τὸν μνηστῆρα … ἐνταῦθα γὰρ εἰσὶν οἱ μάρτυρες, αἱ θεραπαινίδες δηλονότι· οὔκουν ἀπαιτήσει τοὺς μάρτυρας, ἀλλὰ διαβαλεῖ, ὅτι μὴ δεῖ πιστεύειν· «φύσει γὰρ ἐχθρὸν τὸ δοῦλον τοῖς δεσπόταις» ; Plaut.~\textit{Mil.}~475–477 quid propius fuit, / quam ut perirem, si elocutus essem ero? … mussitabis: plus oportet scire servom quam loqui ; \textit{idem} Tosi n°~1978 ; Sen.~\textit{Ep.}~47.4–5 Sic fit, ut isti de domino loquantur, quibus coram domino loqui non licet … Deinde eiusdem arrogantiae proverbium iactatur, totidem hostes esse quot servos ; Tosi n°~239 (Totidem hostes esse quot servos) Totidem hostes esse quot servos … φύσει γὰρ ἐχθρὸν τὸν δοῦλον τοῖς δεσπόταις ; Erasm.~\textit{Adag.}~2.3.31 Quot servos habemus, totidem habemus hostes … Admonet adagium, ne quid fidamus illis, neve ex animo nobis amicos arbitremur, qui nos metu colunt ; Pl.~\textit{Lg.}~757a δοῦλοι γὰρ ἂν καὶ δεσπόται οὐκ ἄν ποτε γένοιντο φίλοι ; Zen.~4.75 Κωρυκαῖος ἠκροάζετο: ταύτης Μένανδρος μέμνηται ἐν τῷ Ἐγχειριδίῳ … καὶ τοῖς πειραταῖς ἐμήνυον}
\lp{340–342 δούλη τινὸς προαγωγὸς οὖσ᾽ … πεμπομένη]}{Eur.~\textit{Hipp.}~645–650 χρῆν δ᾽ ἐς γυναῖκα πρόσπολον μὲν οὐ περᾶν, ἄφθογγα δ᾽ αὐταῖς συγκατοικίζειν δάκη θηρῶν … νῦν δ᾽ αἳ μὲν ἔνδον δρῶσιν αἱ κακαὶ κακὰ βουλεύματ᾽, ἔξω δ᾽ ἐκφέρουσι πρόσπολοι ; Lys.~1.15–16 προσέρχεταί μοί τις πρεσβῦτις ἄνθρωπος, ὑπὸ γυναικὸς ὑποπεμφθεῖσα ἣν ἐκεῖνος ἐμοίχευεν … ἐὰν οὖν λάβῃς τὴν θεράπαιναν τὴν εἰς ἀγορὰν βαδίζουσαν καὶ διακονοῦσαν ὑμῖν καὶ βασανίσῃς, ἅπαντα πεύσῃ ; Lys.~1.8 ἐπιτηρῶν γὰρ τὴν θεράπαιναν τὴν εἰς τὴν ἀγορὰν βαδίζουσαν καὶ λόγους προσφέρων ἀπώλεσεν αὐτήν}
\lp{341 προαγωγὸς οὖσ᾽ ἐνετρύλλισεν τῷ δεσπότῃ]}{\textit{Test.} Suid.~ε~1271 Ἐνετρύλλισε: κατεῖπεν. Ἀριστοφάνης· προαγωγὸς οὖσ᾽ ἐνετρύλλισε τῷ δεσπότῃ ; \textit{idem} Zonar. col.~745~Tittmann (Ἐνετρύλλησε ; οὔ σ᾽ ἐνετρύλλησεν) ; Σ Ar.~\textit{Thesm.}~341 (Dübner, p.~266,~43) κατεῖπε πρὸς τὸν δεσπότην}
\lp{341 προαγωγός]}{\textit{Test.} Hsch.~π~3308 προαγωγός· διδάσκαλος κακῶν. καὶ μαστροπός. †μόνος†. μαυλιστής. ⟨καὶ ἡ ἐπ᾽ αἰσχροῖς ἐπὶ τὸ πορνεῦσαι προάγουσά τινας⟩ (Ar. Thesm. 341) ; \textit{idem} Suid.~π~2307 (om. καὶ μαστροπὸς κτλ.), Phot.~π~1172 Theodoridis (om. καὶ μαστροπὸς κτλ.) ; Ar.~\textit{Thesm.}~558 ὥς τ᾽ αὖ τὰ κρέ᾽ ἐξ Ἀπατουρίων ταῖς μαστροποῖς διδοῦσαι ; Ar.~\textit{Ran.}~1079 οὐ προαγωγοὺς κατέδειξ᾽ οὗτος ; Herod.~1, titre ΠΡΟΚΥΚΛΙΣ Η ΜΑΣΤΡΟΠΟΣ ; Aeschin.~1.184 καὶ τὰς προαγωγοὺς καὶ τοὺς προαγωγοὺς γράφεσθαι κελεύει, κἂν ἁλῶσι, θανάτῳ ζημιοῦν ; Σ Ar.~\textit{Thesm.}~1172 (Dübner, p.~271) Εὐριπίδης ἐν σχήματι προαγωγοῦ γραός ; Hsch.~κ~3626 κορινθιάζειν· μαστροπεύειν, ἑταιρεῖν (Ar. fr. 370 K.–A.) ; Hsch.~κ~4654 Κυπρογενέος προπόλον· προαγωγόν ; Hsch.~μ~370 μαστροπός· δύστροπος, πανοῦργος. ἀπατεών. ὁ τὰς γυναῖκας ἢ ἄνδρας προσκαλῶν καὶ μαυλίζων, ἢ προαγωγός ; Poll.~7.201 (Bethe II, p.~107,~8–10) εἰ δὲ χρὴ καὶ τὰς αἰσχίους πράξεις τέχνας ὀνομάζειν, ἐρεῖς πόρναι πόρνοι, ἑταῖραι ἑταιροῦντες, πορνεῖα οἰκήματα ματρυλεῖα, πορνοβοσκοὶ προαγωγοὶ μαστροποί ; Ar.~\textit{Nub.}~980 αὐτὸς ἑαυτὸν προαγωγεύων τοῖς ὀφθαλμοῖς ἐβάδιζεν ; Ar.~\textit{Vesp.}~1028 ἵνα τὰς Μούσας αἷσιν χρῆται μὴ προαγωγοὺς ἀποφήνῃ ; Pl.~\textit{Tht.}~150a διὰ τὴν ἄδικόν τε καὶ ἄτεχνον συναγωγὴν ἀνδρὸς καὶ γυναικός, ᾗ δὴ προαγωγία ὄνομα ; X.~\textit{Symp.}~4.61 τὴν προαγωγείαν, ἔφη ; Aeschin.~1.14 τὸν τῆς προαγωγείας, τὰ μέγιστα ἐπιτίμια ; Plu.~\textit{Sol.}~23.1 κἂν προαγωγεύῃ, δραχμὰς εἴκοσι ; Σ Ar.~\textit{Nub.}~980 (Dübner, p.~121) προαγωγεύων: μαυλίζων. προαγωγὸς γὰρ μαυλίστρια ; Hsch.~δ~2137 δολάνα· μαστροπός ⟨Λάκωνες⟩ ; Hsch.~κ~786 κάρβις· μαστροπός ; Hsch.~μ~405 ματρυλεῖον· τόπος [τῶν πορνευόντων, τουτέστι πορνεῖον,] ὅπου οἱ μαστροποὶ [ἤτοι μαυλισταί,] διέτριβον ; Hsch.~π~3618 πρόπαις· ὁ προαγωγός, μαστροπός. καὶ μάντις [μαστρόπευσις] ; Synag.~μ~45 (Cunningham, p.~330) μαστροπός· μαυλιστής, πορνοβοσκός ; \textit{idem} Suid.~μ~271 ; Suid.~π~3114 Πυγοστόλος: προαγωγός· μαυλίστρια. Ἡσίοδος· μηδέ σε γυνὴ πυγοστόλος ἐξαπατάτω ; Hsch.~δ~2161 δόλοπα· κατάσκοπον. μαστροπόν ; Hsch.~μ~418 μαυλιστήριον· παρ᾽ Ἱππώνακτι, Λύδιόν τι λεπτὸν νόμισμα (fr. 160 W. = 170 Degani) ; Hsch.~π~3573 πρὸ μὲν λέγοντες· προαγωγοί ; Suid.~μ~270 Μαστροπεύω. Μαστροπεία ; Suid.~μ~272 Μαστρίλλιον· τόπος ἔνθα οἱ μαστροποὶ διέτριβον}
\lp{341 ἐνετρύλλισεν]}{Hsch.~α~8769 Ἀφροδίτη ψίθυρος· … (Cratin. fr. 386 K.–A.) ; Hsch.~τ~1553 τρύζει· γογγύζει, ψιθυρίζει ; \textit{idem} Synag.~τ~274 (Cunningham, p.~475) (om. γογγύζει) ; Hsch.~τ~1554 τρύζητε· ψιθυρίζητε, λογοποιῆτε, γογγύζητε (Ι 311) ; Hsch.~τ~1596 τρώζειν· ψιθυρίζειν. συνουσιάζειν ; Hsch.~ψ~189 ψιθυρίζει· εἰς τὸ οὖς ἠρέμα διαλέγεται ; \textit{idem} Synag.~ψ~18 (Cunningham, p.~519) (ἠρέμα εἰς τὸ οὖς) ; Poll.~5.89 (Bethe I, p.~286,~5–6) καὶ ὄρτυγας τρυλίζειν, καὶ κύκνους ᾄδειν, καὶ τρυγόνας τρύζειν, καὶ περιστερὰς γογγύζειν ; Poll.~7.136 (Bethe II, p.~89,~19–21) στυφοκόμπους δ᾽ αὐτοὺς οἱ κωμῳδοὶ καλοῦσιν, ὡς τὸ ἔκοπτον τοὺς ὄρτυγας καὶ τὸ ἐτρύλιζον καὶ τὸ ἀνεκνάδαλλον τῶν ὀρτυγοκοπικῶν ; Poll.~9.109 (Bethe II, p.~177,~30–32) τοὺς δὲ ἡττηθέντας [ὄρτυγας] ἐμβοήσαντες κατὰ τὸ οὖς αὐτοῖς ἐξιῶντο, λήθην ἐνεργαζόμενοι τῆς τοῦ νενικηκότος φωνῆς· καὶ τὸ ἐμβοᾶν ἐντρυλίζειν ὠνόμαζον ; Harp. s.v.~ψιθυριστὴς Ἑρμῆς ΨΙΘΥΡΙΣΤΗΣ ἙΡΜΗΣ Δημοσθένης ἐν τῷ κατὰ Νεαίρας. ἦν τις Ἀθήνησιν Ἑρμῆς οὕτω καλούμενος ; Hp.~\textit{Int.}~6 οἷον γαστὴρ τρυλλίζει ; Hsch.~δ~1023 διαθρυλλῶ· διαομιλῶ. διαψιθυρίζω ; Hsch.~θ~793 θρύλλοι· ψιθυρισμοί, ὁμιλίαι ; Hsch.~σ~1315 σοίθης· ψίθυρος. ἀλαζών. διάβολος ; Hsch.~τ~1121 τονθορύζει· ἀτάκτως λαλεῖ. γογγύζει. ψιθυρίζει. καὶ τονθρίζει τὰ αὐτά ; Hsch.~ψ~179 ψιδόνες· διάβολοι. ψίθυροι ; Ar.~\textit{Av.}~1297–1299 Μειδίας δ᾽ ἐκεῖ / ὄρτυξ ἐκαλεῖτο· καὶ γὰρ ᾔκειν ὄρτυγι / ὑπὸ στυφοκόπου τὴν κεφαλὴν πεπληγμένῳ ; Hsch.~σ~285 σαυσαρόν· ψίθυρον ; Hsch.~τ~1568 τρυσμός· γογγυσμός ; Hsch.~υ~680 ὑπολαλεῖ· ψιθυρίζει}
\lp{341 ἐνετρύλλισεν τῷ δεσπότῃ]}{Juv.~4.110 Pompeius tenui iugulos aperire susurro}
\lp{341–346 προαγωγός … προδιδοῦσ᾽ ἑταίρᾳ τὸν φίλον]}{Prop.~4.5.1–2 terra tuum spinis obducat, lena, sepulcrum, / et tua, quod non uis, sentiat umbra sitim ; Ov.~\textit{Am.}~1.8.113–114 Di tibi dent nullosque Lares inopemque senectam, / Et longas hiemes perpetuamque sitim! ; Tib.~2.6.53–54 tunc tibi, lena, precor diras: satis anxia uiuas, / mouerit e uotis pars quotacumque deos}
\lp{342 ἀγγελίας ψευδεῖς φέρει]}{Ar.~\textit{Ach.}~114 ἄλλως ἄρ᾽ ἐξαπατώμεθ᾽ ὑπὸ τῶν πρέσβεων ; Eur.~\textit{IT}~461–462 οὐδ᾽ ἀγγελίας ψευδεῖς ἔλακεν βουφορβὸς ἀνήρ ; D.~19.70 λέγε τὴν ἀρὰν καὶ ἀνάγνωθι λαβὼν τὴν ἐκ τοῦ νόμου ταυτηνί. ταῦθ᾽ ὑπὲρ ὑμῶν, ὦ ἄνδρες Ἀθηναῖοι, καθ᾽ ἑκάστην τὴν ἐκκλησίαν ὁ κῆρυξ εὔχεται νόμῳ προστεταγμένα ; Hsch.~ψ~118 ψευδάγγελος· ψευδῆ ἀγγέλ⟨λ⟩ων (Ο 159) ; Ar.~\textit{Av.}~1340 ἔοικεν οὐ ψευδαγγελήσειν ἅγγελος ; Ar.~\textit{Lys.}~1235 ἀγγέλλομεν δ᾽ οὐ ταὐτὰ τῶν αὐτῶν πέρι ; \textit{Il.}~15.159 πάντα τάδ᾽ ἀγγεῖλαι, μὴ δὲ ψευδάγγελος εἶναι ; D.~18.142 ὅτε τοὺς ταλαιπώρους Φωκέας ἐποίησ᾽ ἀπολέσθαι τὰ ψευδῆ δεῦρ᾽ ἀπαγγείλας ; D.~19.69 οἱ δέ, Ἀθηναῖοι, τῆς ἐλευθερωτάτης πόλεως, πρέσβεις ταχθέντες ὑμᾶς, … τούτους ἐξαπατᾶν ὑπέστησαν ; D.~19.76 πέντε γὰρ ἡμέραι γεγόνασι μόναι, ἐν αἷς οὗτος ἀπήγγειλε τὰ ψευδῆ ; App. Prov.~3.81 Μέγα τοι ἄγγελος ἐσθλός: ἐπὶ τῶν ἀγαθὰ ἐπαγγελλόντων ; X.~\textit{Eq.\,mag.}~5.8 φοβεῖν γε μὴν τοὺς πολεμίους καὶ ψευδενέδρας οἷόν τε καὶ ψευδοβοηθείας καὶ ψευδαγγελίας ποιοῦντα ; Arist.~\textit{Po.}~1455a οἷον ἐν τῷ Ὀδυσσεῖ τῷ ψευδαγγέλῳ ; Hsch.~α~377 ἀγγελίας· ἀπαγγελίας (β 92..). σημαίνει δὲ ἀγγελίην αὐτὸν τὸν ἄγγελον (Δ 384..) ; Hsch.~α~378 ἀγγελιαφόρος· πρεσβευτής ; Zen.~2.58 Ἀράβιος ἄγγελος: παροιμία παρὰ τὸ Ἀράβιος αὐλητὴς παραλαμβανομένη ἐπὶ τῶν ἀπαύστως διαλεγομένων ; Diogenian.~3.25 Ἀχνυμένη σκυτάλη: ἐπὶ τῶν λυπηρὰς ἀγγελίας ἀγγελλόντων}
\lp{342–343 ἢ πεμπομένη τις … ἢ μοιχὸς εἴ τις]}{Eur.~\textit{Hec.}~1178–1179 εἴ τις γυναῖκας τῶν πρὶν εἴρηκεν κακῶς ἢ νῦν λέγων ἔστιν τις ἢ μέλλει λέγειν}
\lp{342–343 ψευδεῖς … ψευδῆ]}{Hsch.~ψ~123 ψευδῆ· ἀλλόκοτα. ἀπάται ; Hsch.~ψ~125 ψευδής· ψεύστης (Eur. Or. 1608)}
\lp{343 ἢ μοιχὸς εἴ τις ἐξαπατᾷ ψευδῆ λέγων]}{D.~18.282 καίτοι τίς ὁ τὴν πόλιν ἐξαπατῶν; οὐχ ὁ μὴ λέγων ἃ φρονεῖ; τῷ δ᾽ ὁ κῆρυξ καταρᾶται δικαίως ; D.~23.97 διόπερ καταρᾶται καθ᾽ ἑκάστην ἐκκλησίαν ὁ κῆρυξ, οὐκ εἴ τινες ἐξηπατήθησαν, ἀλλ᾽ εἴ τις ἐξαπατᾷ λέγων ἢ βουλὴν ἢ δῆμον ἢ τὴν ἡλιαίαν ; Din.~1.47 κατάρατος δὲ καθ᾽ ἑκάστην ἐκκλησίαν γιγνόμενος, ἐξεληλεγμένος δῶρα κατὰ τῆς πόλεως εἰληφώς, ἐξηπατηκὼς δὲ καὶ τὸν δῆμον καὶ τὴν βουλὴν παρὰ τὴν ἀράν ; Ar.~\textit{Ran.}~1085–1086 καὶ βωμολόχων δημοπιθήκων ἐξαπατώντων τὸν δῆμον ἀεί}
\lp{343 μοιχός]}{Zen.~4.67 Κέρδος αἰσχύνης ἄμεινον· ἕλκε μοιχὸν εἰς μυχόν ; App. Prov.~5.43 Ὦ Λακιάδαι: ἐπὶ τῶν μοιχῶν. δῆμος γὰρ τῆς Ἀττικῆς αἱ Λακιάδαι, ἐν ᾧ ῥαφανῖδες πολλαί, αἷς ἐχρῶντο κατὰ τῶν ληφθέντων μοιχῶν ἐνυβρίζοντες ; Ar.~\textit{Nub.}~1083–1084 τί δ᾽ ἢν ῥαφανιδωθῇ πιθόμενός σοι τέφρᾳ τε τιλθῇ, / ἕξει τινὰ γνώμην λέγειν τὸ μὴ εὐρύπρωκτος εἶναι ; Ar.~\textit{Lys.}~107 ἀλλ᾽ οὐδὲ μοιχοῦ καταλέλειπται φεψάλυξ ; Ar.~\textit{Lys.}~212–214 οὐκ ἔστιν οὐδεὶς οὔτε μοιχὸς οὔτ᾽ ἀνήρ / … / ὅστις πρὸς ἐμὲ πρόσεισιν ἐστυκώς ; Ar.~\textit{Thesm.}~397 μὴ μοιχὸς ἔνδον ᾖ τις ἀποκεκρυμμένος ; Ar.~\textit{Thesm.}~415–417 σφραγῖδας ἐπιβάλλουσιν ἤδη καὶ μοχλοὺς / τηροῦντες ἡμᾶς, καὶ προσέτι Μολοττικοὺς / τρέφουσι μορμολυκεῖα τοῖς μοιχοῖς κύνας ; Ar.~\textit{Eccl.}~225a μοιχοὺς ἔχουσιν ἔνδον ὥσπερ καὶ πρὸ τοῦ ; Suid.~μ~1359 Μοιχός: ὁ πόρνος. ἐκ τοῦ μὴ καὶ τοῦ οἴχομαι. ὁ ἐπιτηρῶν μὴ οἴχεται ὁ ἀνὴρ τῆς γυναικός, ἵνα εἰσέλθῃ ; Suid.~ρ~55 Ῥαφανιδωθῆναι καὶ τέφρᾳ τιλθῆναι, παρὰ Ἀριστοφάνει. ἐπὶ τῶν ἀσελγῶν. οὕτω γὰρ τοὺς ἁλόντας μοιχοὺς ᾐκίζοντο· ῥαφανῖδας λαμβάνοντες καθίεσαν εἰς τοὺς πρωκτοὺς τούτων καὶ παρατίλλοντες αὐτοὺς θερμὴν τέφραν ἐπέπαττον ; Diogenian.~8.58 Τενέδιος πέλεκυς: ἐπὶ τῶν πικρῶς ἢ μᾶλλον συντόμως ἀποκοπτόντων τὰς ζητήσεις … Νόμον γάρ τις ἔταξε παρ᾽ αὐτοῖς, τὸν καταλαβόντα μοιχὸν πελέκει κόπτειν ; Ar.~\textit{Nub.}~1079 μοιχὸς γὰρ ἢν τύχῃς ἁλούς, τάδ᾽ ἀντερεῖς πρὸς αὐτόν ; Ar.~\textit{Eccl.}~522 οὔτοι παρὰ τοῦ μοιχοῦ γε φήσεις ; Ar.~fr.~191~K.-A. (187~Kock) πάσαις γυναιξὶν ἐξ ἑνός γε του ⟨τρόπου⟩ / ὥσπερ παροψὶς μοιχὸς ἐσκευασμένος ; Hsch.~μ~1558 μοιχάγρια· τὰ τῆς μοιχείας ἀγρεύματα· ὁ γὰρ ληφθεὶς ἐπὶ μοιχείᾳ ζημιοῦται (θ 332) ; Hsch.~μ~1559 μοιχοῖς· πόρνοις ; Hsch.~μ~1560 μοιχοτύπη· ἡ ὑπὸ μοιχῶν τυπτομένη (com. ad. fr. 389 K.–A.) ; Suid.~μ~1354 Μοιχάγρια: τὰ ἀντὶ τῆς μοιχείας ἀγρεύματα, ἐπὶ τῷ ληφθῆναι μοιχὸν ἀποτινόμενα ; Suid.~μ~1357 μοιχεία καὶ φαρμακεία ταυτόν ἐστιν. καὶ γὰρ ἀμφότερα λάθρα γίνεται. λέγονται δὲ καὶ ἐπιβουλαί· καὶ οὔτε ἡ μοιχευομένη οὔτε ἡ φαρμακεύουσα φιλεῖ τὸν ἄνδρα ; Erasm.~\textit{Adag.}~1.9.29 Tenedia bipennis. Τενέδιος πέλεκυς. i. e. Tenedia securis. Conveniet aut in eos, qui nimium severe judicant ac puniunt … aut quod Tenedius quidam Rex legem tulit in adulteros, ut utriusque corpus securi concideretur}
\lp{343 ψευδῆ λέγων]}{Eur.~\textit{IA}~1006 ψευδῆ λέγων δὲ καὶ μάτην ἐγκερτομῶν}
\lp{343 ἐξαπατᾷ]}{App. Prov.~2.71 Ἐξηπάτησεν ἡ χάραξ τὴν ἄμπελον: ὅταν τὸ σωζόμενον ὑπὸ τοῦ σώζοντος ἐξαπατηθῇ ; Hes.~\textit{Op.}~373 μὴ δὲ γυνή σε νόον πυγοστόλος ἐξαπατάτω ; Hsch.~α~5835 ἀπατεύων· ψευδόμενος ; Hsch.~ε~3580 ἐξαπατᾷ· χλευάζει, δελεάζει, φενακίζει (Ar. Av. 521) ; Hsch.~η~675 ἠπεροπεύειν· ἐξαπατᾶν, καὶ ψεύδεσθαι, ἤτοι παραλογίζεσθαι (Γ 399..) ; App. Prov.~4.45 Οὐκ ἀπείρῳ προςέβαλες: ἐπὶ τῶν οἰομένων μὲν ἐξαπατᾶν, φωραθέντων δέ ; Hsch.~σ~330 σείεις· ἐξαπατᾷς. Ἀττικοί (cf. Antiph. 6, 43)}
\lp{343–344 ἢ μοιχὸς … ἐξαπατᾷ ψευδῆ λέγων καὶ μὴ δίδωσιν ἃν ὑπόσχηταί ποτε]}{Philonid.~fr.~7~Kock ὅρκους δὲ μοιχῶν εἰς τέφραν ἐγὼ γράφω ; Call.~\textit{Epigr.}~25.3–4~Pf. ὤμοσεν· ἀλλὰ λέγουσιν ἀληθέα τοὺς ἐν ἔρωτι ὅρκους μὴ δύνειν οὔατ᾽ ἐς ἀθανάτων ; Pl.~\textit{Phlb.}~65c καὶ ἐν ταῖς ἡδοναῖς ταῖς περὶ τἀφροδίσια, αἳ δὴ μέγισται δοκοῦσιν εἶναι, καὶ τὸ ἐπιορκεῖν συγγνώμην εἴληφε παρὰ θεῶν ; Pl.~\textit{Smp.}~183b ὅτι καὶ ὀμνύντι μόνῳ συγγνώμη παρὰ θεῶν ἐκβάντι τῶν ὅρκων—ἀφροδίσιον γὰρ ὅρκον οὔ φασιν εἶναι ; Hsch.~α~8771 ἀφροδίσιος ὅρκος· παροιμία, ἣν καὶ ἀναγράφουσιν· Ἀφροδίσιος ὅρκος οὐ δάκνει· πρῶτος δὲ Ἡσίοδος (fr. 124 M.–W.) ἔπλασε τὰ περὶ τὸν Δία καὶ τὴν Ἰώ. «ὤμοσεν, ἀλλὰ ⟨λέγουσιν⟩ …» (Callim. ep. 25, 3) ; Diogenian.~3.37 Ἀφροδίσιος ὅρκος οὐκ ἐμποίνιμος: ἐπὶ τῶν ἐρωτικῶν ; \textit{idem} Greg. Cypr.~1.2 (om. ἐπὶ τῶν ἐρωτικῶν), Suid.~α~4652 (καὶ παροιμία· … ἐπὶ τῶν δι᾽ ἔρωτα ὀμνυόντων πολλάκις καὶ ἐπιορκούντων), Tosi n°~332 (Iuppiter ex alto periuria ridet amantum) (om. ἐπὶ τῶν ἐρωτικῶν) ; App. Prov.~4.33 Ὅρκος Ἀφροδίσιος συγγινώσκεται: παρόσον οἱ ἐρῶντες εὐχερῶς ὀμνύουσι ; Erasm.~\textit{Adag.}~2.4.90 Venereum jusjurandum. Ἀφροδίσιος ὅρκος οὐκ ἐμποίνιμος. i. e. Venereum jusjurandum non punitur. Sensus est, amantium jusjurandum irritum esse: neque ad Deos quicquam pertinere, tanquam jocosum ; Xenarch.~\textit{Pentathl.} ap. Ath.~10.441e ὅρκον δ᾽ ἐγὼ γυναικὸς εἰς οἶνον γράφω ; Plaut.~\textit{Cist.}~98–103 At ille conceptis iuravit verbis apud matrem meam, / me uxorem ducturum esse: ei nunc alia ducendast domum … eum uxorem ducturum esse aliam. — Nihil amori iniuriumst ; Soph.~fr.~811~R. (742~N.²) ὅρκους ἐγὼ γυναικὸς εἰς ὕδωρ γράφω ; \textit{idem} Tosi n°~337 (Mulier cupido quod dicit amanti … rapida scribere oportet aqua) ; Hor.~\textit{Carm.}~2.8.1–2,~13–14 Vlla si iuris tibi peierati / poena, Barine, nocuisset umquam … ridet hoc, inquam, Venus ipsa, rident / simplices Nymphae ; Ov.~\textit{Am.}~1.8.85–86 Nec, siquem falles, tu periurare timeto — / Commodat in lusus numina surda Venus ; Tib.~1.4.21–22 Nec iurare time: Veneris periuria venti / Inrita per terras et freta summa ferunt ; Apollod.~2.1.3 ἀπωμόσατο δὲ ταύτῃ μὴ συνελθεῖν· διό φησιν Ἡσίοδος οὐκ ἐπισπᾶσθαι τὴν ἀπὸ τῶν θεῶν ὀργὴν τοὺς γινομένους ὅρκους ὑπὲρ ἔρωτος ; Σ \textit{Pl.\,Symp.}~183b (ap. Leutsch–Schneidewin, CPG I, p.~221, adn. ad Diogenian.~3.37) ἐπὶ τῶν δι᾽ ἔρωτα ὀμνυόντων πολλάκις καὶ ἐπιορκούντων· μέμνηται δὲ ταύτης καὶ Ἡσίοδος, λέγων … Ἐκ τοῦδ᾽ ὅρκον ἔθηκεν ἀμείνονα ἀνθρώποισι / Νοσφιδίων ἔργων περὶ Κύπριδος ; Hsch.~τ~287 ταχυβάμονας ὅρκους· ἀντὶ τοῦ βραδεῖς, βέλτιον δὲ τοὺς ἐν ἔρωτι γινομένους· ἀβλαβεῖς γὰρ οἱ ἐπιορκήσαντες (trag. ad. fr. 333 a K.–S.) ; Publil. Syr. A~37 Meyer Amantis ius iurandum poenam non habet ; Suid.~ει~314 Εἰς τέφραν γράφειν: ἐπὶ τῶν ἀδυνάτων. ἐγὼ δὲ τοὺς ὅρκους αὐτοῦ, τὸ τῆς παροιμίας, οἶμαι δεῖν εἰς τέφραν γράφειν ; Suid.~ε~3207 Ἐς τέφραν γράφεις: ἐπὶ τῶν μάτην πονούντων, καὶ ἀδυνάτων}
\lp{343–344 ἐξαπατᾷ … ἃν ὑπόσχηταί ποτε]}{D.~20.100 ἔστι δὲ δήπου νόμος ὑμῖν, ἐάν τις ὑποσχόμενός τι τὸν δῆμον ἢ τὴν βουλὴν ἢ δικαστήριον ἐξαπατήσῃ, τὰ ἔσχατα πάσχειν ; D.~20.135 ἔστιν ὑμῖν νόμος ἀρχαῖος, τῶν καλῶς δοκούντων ἔχειν, ἄν τις ὑποσχόμενός τι τὸν δῆμον ἐξαπατήσῃ, κρίνειν, κἂν ἁλῷ, θανάτῳ ζημιοῦν ; D.~49.67 ὃς οὖν οὔτε ὑμᾶς ᾐσχύνθη ἐξαπατῆσαι ὑποσχόμενος, νόμων ὄντων, ἐάν τις τὸν δῆμον ὑποσχόμενος ἐξαπατήσῃ, εἰσαγγελίαν εἶναι περὶ αὐτοῦ ; Arist.~\textit{Ath.}~43.5 καὶ συκοφαντῶν προβολὰς τῶν Ἀθηναίων καὶ τῶν μετοίκων μέχρι τριῶν ἑκατέρων, κἄν τις ὑποσχόμενός τι μὴ ποιήσῃ τῷ δήμῳ ; Arist.~\textit{Ath.}~28.3 πρῶτος ὑποσχόμενος ἐπιθήσειν}
\lp{344 καὶ μὴ δίδωσιν ἃν ὑπόσχηταί ποτε]}{Ov.~\textit{Ars}~1.443–444 Promittas facito: quid enim promittere laedit? / Pollicitis dives quilibet esse potest ; Ov.~\textit{Ars}~1.631–634 Nec timide promitte: trahunt promissa puellas; / Pollicito testes quoslibet adde deos. / Iuppiter ex alto periuria ridet amantum, / Et iubet Aeolios inrita ferre notos ; Diogenian.~2.1 Αἱ Χάρητος ὑποσχέσεις: ἐπὶ τῶν προχείρως ἐπαγγελλομένων ; \textit{idem} App. Prov.~5.24 (om. Αἱ ; εὐχερῶς ; add. Χάρης γὰρ ὁ στρατηγὸς Ἀθηναίων εὐχερής τε καὶ εὐκολώτατος πρὸς τὰς ὑποσχέσεις ἐγένετο) ; Diogenian.~6.12 Λευκὴν μᾶζαν φυρῶ σοι: ἐπὶ τῶν μεγάλως ὑπισχνουμένων ; App. Prov.~3.20 Θετταλικὸν σόφισμα: ἐπὶ τῶν πανούργων καὶ εὐμεθόδων ἔργων. Λέγεται δὲ καὶ ἐπὶ τῶν ἀεὶ ἀναβαλλομένων τὸ παρ᾽ αὐτῶν ὑπεσχημένον, δι᾽ οἰκεῖον κέρδος ; \textit{idem} Zen.~4.29 (Θετταλῶν σόφισμα· παροιμία ἐπὶ τῶν σοφιζομένων λεγομένη) ; Erasm.~\textit{Adag.}~2.6.84 Charetis pollicitationes. Αἱ Χάρητος ὑποσχέσεις. i. e. Charetis pollicitationes. Dici solitum, ubi quis facile ac benigne promittit … Natum adagium à Charete quodam Atheniensium duce, qui quidvis cuivis pollicebatur, at non praestabat item quod receperat ; Catull.~70.3–4 sed mulier cupido quod dicit amanti / in vento et rapida scribere oportet aqua ; Sall.~\textit{Cat.}~23.3 erat ei cum Fulvia, muliere nobili, stupri vetus consuetudo. quoi cum minus gratus esset, quia inopia minus largiri poterat, repente glorians maria montisque polliceri coepit ; \textit{idem} Tosi n°~1505 ; Hsch.~υ~553 ὑπισχνεῖται· ὁμολογεῖ, συντίθεται, ἐπαγγέλλεται ; Greg. Cypr.~2.68 Θάσος ἀγαθῶν: ἐπὶ τῶν εὐδαίμονα καὶ λαμπρὸν ἀποδεῖξαί τινα ἐπαγγελλομένων ; App. Prov.~2.9b Ἐγώ τοι πάντα ποιήσω θέρος: ἐπὶ τῶν ἀγαθὰ ὑπισχνουμένων δώσειν ; Tosi n°~442 Facile ex amico inimicum facies, cui promissa non reddas ; Erasm.~\textit{Adag.}~1.3.10 Thessalorum commentum. Θετταλῶν σόφισμα, i. e. Thessalorum commentum, varias ad res accommodari scribit Suidas … tum ad hos qui dolis quempiam circumveniunt … votum nuncupat, se centum hominum capita Deo sacrificaturum … ut postea quotannis hecatomben ei Deo polliceri sint soliti: neque reddere tamen ; Ar.~\textit{Thesm.}~217 ἢ μὴ ᾽πιδοῦν᾽ ἐμαυτὸν ὤφελόν ποτε ; Ar.~\textit{Plu.}~865 ὑποσχόμενος οὗτος ποιήσειν εὐθέως ; Hsch.~υ~766 ὑποσχέσεις· ἐπαγγελίαι ; Hsch.~υ~767 ὑπόσχωμαι· ὑποβάλωμαι, ἐπαγγέλλωμαι (Χ 114)}
\lp{344–345 καὶ μὴ δίδωσιν … ἢ δῶρά τις δίδωσι]}{Ar.~\textit{Plu.}~153–159 καὶ τούς γε παῖδάς φασι ταὐτὸ τοῦτο δρᾶν / οὐ τῶν ἐραστῶν ἀλλὰ τἀργυρίου χάριν … αἰσχυνόμενοι γὰρ ἀργύριον αἰτεῖν ἴσως / ὀνόματι περιπέττουσι τὴν μοχθηρίαν}
\lp{345 ἢ δῶρά τις δίδωσι μοιχῷ γραῦς γυνή]}{Ar.~\textit{Plu.}~959–961 ἆρ᾽ ὦ φίλοι γέροντες ἐπὶ τὴν οἰκίαν / ἀφίγμεθ᾽ ὄντως τοῦ νέου τούτου θεοῦ ; Ar.~\textit{Plu.}~975–979 ἦν μοί τι μειράκιον φίλον, πενιχρὸν μέν, ἄλλως δ᾽ εὐπρόσωπον καὶ καλὸν καὶ χρηστόν· εἰ γάρ του δεηθείην ἐγώ, ἅπαντ᾽ ἐποίει κοσμίως μοι καὶ καλῶς· ἐγὼ δ᾽ ἐκείνῳ πάντα ταῦθ᾽ ὑπηρέτουν ; Ar.~\textit{Plu.}~1024 γραὸς καπρώσης τἀφόδια κατεσθίειν ; Plu.~\textit{Sol.}~20.5 καὶ νέον ἐν δωματίῳ πλουσίας πρεσβύτιδος, ὥσπερ οἱ πέρδικες, ἀπὸ συνουσίας παχυνόμενον ; Σ Ar.~\textit{Thesm.}~345 (Dübner, p.~266,~44) ἵνα αὐτῇ συγγένηται ; Ar.~\textit{Eccl.}~877–881 τί ποθ᾽ ἅνδρες οὐχ ἥκουσιν; ὥρα δ᾽ ἦν πάλαι· / ἐγὼ δὲ καταπεπλασμένη ψιμυθίῳ / ἕστηκα … πῶς ἂν περιλάβοιμ᾽ αὐτῶν τινὰ ; Ar.~\textit{Eccl.}~893 εἴ τις ἀγαθὸν βούλεται παθεῖν}
\lp{345 γραῦς γυνή]}{Zen.~2.96 Γραῦς βακχεύει· ἐπὶ τῶν παρ᾽ ὥραν τι διαπραττομένων. Ταῖς γὰρ νέαις πρέπει τὸ χορεύειν ; \textit{idem} Diogenian.~3.74 (ποιούντων ; om. Ταῖς … χορεύειν ; add. Καὶ Ἀριστοφάνης, Γραῦς καπρῶσα καὶ βακχεύουσα), Greg. Cypr.~1.99 (om. Ταῖς … χορεύειν), Suid.~γ~432 (Γραῦς χορεύει) ; Zen.~6.42 Χαῖρε φίλον φῶς: γραῦς θέλουσα ἀκολασταίνειν γυμνὴ, ἵνα μὴ τὴν ῥάκωσιν τοῦ σώματος ἐλέγξῃ, τὸν λύχνον ἀποσβέσασα εἶπε ; Diogenian.~4.10 Γραῦς ἀναθυᾷ: ἀντὶ τοῦ, καπρᾷ. Ἐπὶ τῶν παρ᾽ ἡλικίαν γυναικιζομένων. Ὡς καὶ τὸ, Γραῦς ἀναβακχεύει ; App. Prov.~3.2 Ἡ γραῦς τὰς ἀγαθίδας ** ἐκπίνουσα, ταῦτα μὲν ἀγαθίδες Ἐπὶ τῶν ἄγαν μεθύσων ; Erasm.~\textit{Adag.}~1.9.8 Anus bacchatur. Γραῦς βακχεύει. i. e. Anus bacchatur. In eos dici solitum, qui praeter aetatem indecore lasciviunt … Quandoquidem insaniae species est, aut vetulam amare, aut potare, aut choros ducere ; Erasm.~\textit{Adag.}~1.9.9 Anus hircissans. Γραῦς καπρῶσα. i. e. Anus hircissans. De anu adhuc intempestiva libidine prurienti … Aristophanes in Pluto … De juvene, cui anus libidinosa omnia suppeditabat, quo vicissim ab illo voluptatem auferret ; Eur.~\textit{Hec.}~323 γραῖαι γυναῖκες ἠδὲ πρεσβῦται σέθεν ; Eur.~\textit{Tro.}~490 δούλη γυνὴ γραῦς Ἑλλάδ᾽ εἰσαφίξομαι ; Archil.~fr.~31~Bergk⁴ (ap. Plu.~\textit{Per.}~28.5) οὐκ ἂν μύροισι γραῦς ἐοῦσ᾽ ἠλείφεο ; Hor.~\textit{Carm.}~3.15.1–4,~13–16 Uxor pauperis Ibyci, / tandem nequitiae fige modum tuae … te lanae prope nobilem / tonsae Luceriam, non citharae decent / nec flos purpureus rosae / nec poti, vetulam, faece tenus cadi ; Hor.~\textit{Epod.}~8.1–2 Rogare longo putidam te saeculo, / viris quid enervet meas ; D.~19.283 τὴν μητέρα δεῖν τὴν Τιμάρχου, γραῦν γυναῖκα ; D.~47.81 γραῦν γυναῖκα ; Hsch.~γ~908 γραῦς Ἔριφος· Ἀπολλόδωρός φησι, ὅτι ἔστι τις παροιμιώδης λεγομένη Ἐριφία γραῦς· εἶναι δὲ αὐτὴν οἷον γραῦν ἔριφον, τὴν ἐν παρθενίᾳ καταγεγηρακυῖαν ; Publil. Syr. A~30 Meyer Anus cum ludit, morti delicias facit ; Erasm.~\textit{Adag.}~2.8.11 Anus saltat. Γραῦς χορεύει. i. e. Anus saltat. In eum convenit, qui indecore quippiam facit … Horatius in anum intempestivae libidinis: Te lanae prope nobilem / Tonsae Luceriam, non citharae decent ; Ar.~\textit{Thesm.}~505 τὸ δ᾽ εἰσέφερε γραῦς ἐν χύτρᾳ τὸ παιδίον ; Ar.~\textit{Thesm.}~512 εἶθ᾽ ἡ μιαρὰ γραῦς, ἣ ᾽φερεν τὸ παιδίον ; Ar.~\textit{Thesm.}~896 ξένη τίς ἡ γραῦς ἡ κακορροθοῦσά σε ; Ar.~\textit{Thesm.}~1073 ἀπολεῖς μ᾽ ὦ γραῦ στωμυλλομένη ; Hsch.~γ~907 γραῦς· τὸ συναγόμενον καὶ ἐπιπηγνύμενον ἐπάνω, ὅταν γάλα ἕψηται, ἤ τι τοιοῦτον ὑγρόν (Ar. Plut. 1206). καὶ ἡ ἐν τοῖς χείλεσι τῶν ποτηρίων γραμμή. καὶ βόλος τις ἀστραγάλων ; Suid.~γ~428 Γραῦς: καὶ ἡ ἄνθρωπος, καὶ ὁ ἀνώτατος ἀφρὸς ἐν τῷ ὑπερζεῖν τὰ ὄσπρια γινόμενος. Ἀριστοφάνης Πλούτῳ ; Diogenian.~3.97 Γραῦς ἀνακροτήσασα πολὺν κονιορτὸν ἐγείρει: ἐπὶ τῶν δι᾽ ἐμπειρίαν πολλὰ πράγματα κινούντων ; \textit{idem} Suid.~γ~429, Erasm.~\textit{Adag.}~2.8.12 (praem. Anus subsultans multum excitat pulveris ; om. ἐπὶ τῶν … κινούντων) ; Publil. Syr. A~29 Meyer Amare iuueni fructus est, crimen seni ; Suid.~γ~431 Γραῦς σέριφος: ἡ ἐν παρθενίᾳ γεγηρακυῖα}
\lp{345 ἢ δῶρά τις δίδωσι]}{D.~21.113 ἐάν τις Ἀθηναίων λαμβάνῃ παρά τινος, ἢ αὐτὸς διδῷ ἑτέρῳ ; Hsch.~δ~2735 δωροδοκεῖ· δίδωσι δῶρα}
\lp{345–346 δῶρά τις δίδωσι … δέχεται]}{Hsch.~δ~2737 δωροδοκία· τὸ λαβεῖν ἢ δοῦναι δῶρα ; Ov.~\textit{Ars}~3.653 Munera, crede mihi, capiunt hominesque deosque ; Suid.~δ~1469 Δωροδοκήσαντα: οὐ μόνον δῶρα διδόναι καὶ ἀναπείθειν, ἀλλὰ καὶ τὸ λαμβάνειν}
\lp{345–346 γραῦς γυνή … δέχεται προδιδοῦσ᾽ ἑταίρᾳ τὸν φίλον]}{Phot. s.v.~ματρυλεῖον (Naber I, p.~408) ματρυλεῖον· τόπος ἐν ᾧ γρᾶες μαστροποὶ διατρίβουσιν ἑταίρας ἔχουσαι καὶ δέχονται τοὺς βουλομένους· ἐκάλουν δὲ αὐτὰς καὶ μητέρας εὐφήμως ; Harp. s.v.~ματρυλεῖον Ἡρακλέων καὶ Δίδυμος τόπον τινά φασιν εἶναι ἐν ᾧ γρᾶες διατρίβουσαι δέχονται τοὺς βουλομένους καταμεθυσθῆναι· Μένανδρος Ἐπιτρέπουσιν· οὐκ οἰμώξεται καταφθαρεὶς ἐν ματρυλείῳ τὸν βίον}
\lp{346 ἢ καὶ δέχεται προδιδοῦσ᾽ ἑταίρᾳ τὸν φίλον]}{D.~46.26 ἐπὶ δωροδοκίᾳ χρήματα διδοὺς ἢ δεχόμενος ; Din.~2.16 πρῶτον μὲν καθ᾽ ἑκάστην ⟨ἐκκλησίαν⟩ δημοσίᾳ κατὰ τῶν πονηρῶν ἀρὰς ποιούμενοι, εἴ τις δῶρα λαμβάνων μετὰ ταῦτα λέγει καὶ γιγνώσκει περὶ τῶν πραγμάτων, ἐξώλη τοῦτον εἶναι ; Hyp.~4.8 ἢ ἐάν τις πόλιν τινὰ προδῷ ἢ ναῦς ἢ πεζὴν ἢ ναυτικὴν στρατιάν, ἢ ῥήτωρ ὢν μὴ λέγῃ τὰ ἄριστα τῷ δήμῳ τῷ Ἀθηναίων χρήματα λαμβάνων ; Σ Ar.~\textit{Thesm.}~346 (Dübner, p.~266,~45–50) ἢ καὶ δέχεται δῶρα παρ᾽ ἑταίρας ἐπὶ τῷ προδιδόναι καὶ συγχωρῆσαι αὐτῇ τὸν φίλον. γράφεται ἔν τισιν ἑταίρα καὶ ἔστιν ἕτερος νοῦς. Ἄλλως. ἐὰν μὲν ἄνευ τοῦ ι γρ. ἑταίρα, οὕτως· ἢ δέχεται δῶρα ὑπὲρ τοῦ ἑταίρα οὖσα προδιδόναι τὸν φίλον. ἐὰν δὲ σὺν τῷ ι, προδιδοῦσάν τινα ἑταίραν τὸν ἑαυτῆς φίλον}
\lp{346 ἢ καί]}{Eur.~\textit{Med.}~42 ἢ καὶ τύραννον τόν τε γήμαντα κτάνῃ}
\lp{346 δέχεται]}{Pl.~\textit{R.}~390e οὐ μὲν δὴ δωροδόκους γε ἐατέον εἶναι τοὺς ἄνδρας οὐδὲ φιλοχρημάτους … οὐδ᾽ ᾀστέον αὐτοῖς ὅτι— δῶρα θεοὺς πείθει, δῶρ᾽ αἰδοίους βασιλῆας ; \textit{idem} Diogenian.~4.21 (καὶ αἰδοίους ; add. ἐπὶ τῶν διὰ δῶρα τὰς δίκας ἀντιστρεφόντων), Greg. Cypr.~2.18 (καὶ αἰδοίους), Suid.~δ~1451 (καὶ αἰδοίους ; add. οἱ μὲν Ἡσιόδειον οἴονται τὸν στίχον, εἴρηται δὲ καὶ Πλάτωνος ἐν τρίτῃ Πολιτείᾳ) ; Greg. Cypr.~1.95 Βοῦς ἐπὶ γλώσσης: ἐπὶ τῶν δωροδοκουμένων· διὰ τὸ τῶν Ἀθηναίων νόμισμα βοῦν ἔχειν ; Ar.~\textit{Eq.}~66 αἰτεῖ ταράττει δωροδοκεῖ λέγων τάδε ; Eur.~\textit{Med.}~964–965 πείθειν δῶρα καὶ θεοὺς λόγος· / χρυσὸς δὲ κρείσσων μυρίων λόγων βροτοῖς ; Hsch.~δ~2740 δωροδόκος· δεχόμενος δῶρα (Dem. 19, 201. Aeschin. 3, 244) ; \textit{idem} Synag.~δ~440 (Cunningham, p.~183) (δῶρα δεχόμενος) ; Hsch.~μ~417 μαῦλις· μάχαιρα. καὶ ἡ μίσθιον ποιοῦσα ; Zen.~1.26 Αἴξ οὐρανία: Κρατῖνος φησὶ καθάπερ Διὸς αἶγα Ἀμαλθείαν, οὕτω καὶ τῶν δωροδοκούντων αἶγα οὐρανίαν ; Zen.~5.2 Λύκου δεκάς: παροιμιῶδες. Ἐπειδὴ Λύκος ὁ ἥρως πρὸς τοῖς ἐν Ἀθήναις δικαστηρίοις ἵδρυτο, τοῦ θηρίου τὴν μορφὴν ἔχων, ἔνθα οἱ δωροδοκοῦντες καὶ συκοφάνται κατὰ δέκα γενόμενοι συνεστρέφοντο ; Diogenian.~2.18 Ἀλώπηξ οὐ δωροδοκεῖται: ἐπὶ τῶν οὐ ῥᾳδίως ἁλισκομένων ; Erasm.~\textit{Adag.}~1.3.18 Muneribus vel Dii capiuntur. Quod scripsit Ovidius: Munera crede mihi capiunt hominesque deosque … Δῶρα θεοὺς πείθει, καὶ αἰδοίους βασιλῆας. i. e. Cum divis flectunt venerandos munera Reges ; Erasm.~\textit{Adag.}~1.7.18 Bos in lingua. Βοῦς ἐπὶ γλώσσης. i. e. Bos in lingua. In eos, qui non audent libere, quod sentiant, dicere … Vel hinc quod Atheniensium nomisma quondam bovis obtinuit figuram … qui pecunia corrupti loqui non auderent, bovem in lingua dicebantur habere ; Ar.~\textit{Av.}~510 ἐπὶ τῶν σκήπτρων ἐκάθητ᾽ ὄρνις μετέχων ὅ τι δωροδοκοίη}
\lp{346 προδιδοῦσ᾽ ἑταίρᾳ τὸν φίλον]}{Ar.~\textit{Thesm.}~1172–1173 ἐμὸν ἔργον ἐστίν· καὶ σὸν ὦλάφιον ἅ σοι / καθ᾽ ὁδὸν ἔφραζον ταῦτα μεμνῆσθαι ποιεῖν ; Epicr.~fr.~9~Kock (Ael.~\textit{NA}~12.10) τελέως δέ μ᾽ ὑπῆλθεν ἡ κατάρατος μαστροπὸς ἐπομνύουσα ; Eur.~\textit{Hipp.}~591 ὤμοι ἐγὼ κακῶν· προδέδοσαι, φίλα ; Diogenian.~5.65 Κρωβύλου ζεῦγος: ἐπὶ τῶν ὑπερβαλλούσῃ πονηρίᾳ κεχρημένων. Ἀπὸ Κρωβύλου τινὸς πορνοβόσκου, δύο ἔχοντος ἑταίρας ; Ar.~\textit{Ran.}~1079 οὐ προαγωγοὺς κατέδειξ᾽ οὗτος ; Eur.~\textit{Med.}~17 προδοὺς γὰρ αὑτοῦ τέκνα δεσπότιν τ᾽ ἐμὴν ; Ov.~\textit{Am.}~1.8.2 est quaedam nomine Dipsas anus ; Ov.~\textit{Am.}~3.5.40 Ingenium dominae lena movebat anus ; Zen.~1.3 Ἀγαθὰ Κιλλικῶν: λείπει τὸ ἔχει. Προδότης γὰρ γέγονεν οὗτος ὁ Κιλλικῶν, Μιλήσιος τὸ γένος, ὃς προδοὺς Μίλητον εὐπόρησεν ; Diogenian.~8.29 Τὸ Φρυνίχου πάλαισμα: κατὰ τῶν πανούργως σοφιζομένων. Προδότης γὰρ γέγονε τῶν Λακεδαιμονίων ; Eur.~fr.~655~N.² οὐκ ἂν προδοίην καίπερ ἄψυχον φίλον ; Hsch.~π~3388 πρόδοσιν· προδομάτιον (Hermipp. fr. 78 K.–A.) ; Suid.~π~2370 Προδότις· θηλυκὸν δὲ διὰ τοῦ ι. ὦ προδότι καὶ παράγωγε καὶ μύραινα σύ ; Diogenian.~5.74 Κακὴ πρόδοσις: ἐπὶ τῶν μετὰ τὴν μίξιν τὸν συγκείμενον μισθὸν κομιζομένων}
\lp{346 ἑταίρᾳ]}{Ar.~\textit{Eccl.}~718–720 ἔπειτα τὰς πόρνας καταπαῦσαι βούλομαι / ἁπαξαπάσας … ἵνα τῶν νέων ἔχωσιν αὗται τὰς ἀκμάς ; Ar.~\textit{Plu.}~149–150 καὶ τάς γ᾽ ἑταίρας φασὶ τὰς Κορινθίας, ὅταν μὲν αὐτάς τις πένης πειρῶν τύχῃ ; Hsch.~ε~6480 ἑταίρα· φίλη (Ι 2). πόρνη. καὶ ἑταιρεύουσα ; \textit{idem} Synag.~ε~885 (Cunningham, p.~235) (om. καὶ ἑταιρεύουσα) ; Ar.~\textit{Lys.}~701 τοῖσι παισὶ τὴν ἑταίραν ἐκάλεσ᾽ ἐκ τῶν γειτόνων ; Ar.~\textit{Eccl.}~528 γυνή μέ τις νύκτωρ ἑταίρα καὶ φίλη ; Σ Ar.~\textit{Lys.}~701 (Dübner, p.~257) ὅτι τὴν γυναῖκα ἑταίραν εἶπεν ὡς ἑταῖρον τὸν συνήθη καὶ φίλον ; Hsch.~ε~6493 ἑταῖρος· ὁ φίλος}
\lp{346 τὸν φίλον]}{Ar.~\textit{Thesm.}~479 ὁ δ᾽ ἀνὴρ παρ᾽ ἐμοὶ καθηῦδεν· ἦν δέ μοι φίλος ; X.~\textit{Mem.}~3.11.4 ἐάν τις, ἔφη, φίλος μοι γενόμενος εὖ ποιεῖν ἐθέλῃ, οὗτός μοι βίος ἐστί ; Ar.~\textit{Eccl.}~931–932 ᾄδω πρὸς ἐμαυτὴν Ἐπιγένει τὠμῷ φίλῳ. / σοὶ γὰρ φίλος τίς ἐστιν ἄλλος ἢ Γέρης}
\lp{347–348 κεἴ τις κάπηλος ἢ καπηλὶς τοῦ χοὸς ἢ τῶν κοτυλῶν τὸ νόμισμα διαλυμαίνεται]}{\textit{Test.} Suid.~χ~394 Χοός: διξέστου· χοῦς γὰρ δύο ξέσται, χοεὺς δὲ ἕξ. Ἀριστοφάνης· εἴ τις κάπηλος ἢ καπηλὶς τοῦ χοὸς ἢ τῶν κοτυλῶν τὸ νόμισμα διαλυμαίνεται ; Ar.~\textit{Plu.}~435–436 ἆρ᾽ ἐστὶν ἡ καπηλὶς ἡκ τῶν γειτόνων, ἣ ταῖς κοτύλαις ἀεί με διαλυμαίνεται}
\lp{347 κεἴ τις κάπηλος ἢ καπηλίς]}{Ar.~\textit{Pax}~447–448 κεἴ τις δορυξὸς ἢ κάπηλος ἀσπίδων, ἵν᾽ ἐμπολᾷ βέλτιον, ἐπιθυμεῖ μαχῶν ; Pl. Com. fr.~174~Kock (Ath.~10.441e) εἶεν, γυναῖκες, ὡς ὑμῖν πάλαι / οἶνον γενέσθαι τὴν ἄνοιαν εὔχομαι. / ὑμῖν γὰρ οὐδέν, καθάπερ ἡ παροιμία, / ἐν τῷ καπήλῳ νοῦς ἐνεῖναί μοι δοκεῖ ; Antiph.~fr.~24~Kock (Ath.~10.441b) γείτων ἐστί τις κάπηλος· οὗτος εὐθὺς ὅταν ἔλθω ποτὲ διψῶσα, μόνος οἶδ᾽ ὥς γ᾽ ἐμοὶ κεράννυται ; Com. adesp.~561~Kock (Macar.~8.33) Τί δ᾽ ἂν κάπηλος παρὰ καπηλίδος λάβοι: ὅτι οὐδὲν ἢ οἶνον ; Nicostr. Com. fr.~22~Kock (Ath.~15.700b) ὁ κάπηλος γὰρ οὑκ τῶν γειτόνων ἄν τ᾽ οἶνον ἄν τε φανὸν ἀποδῶταί τινι ἄν τ᾽ ὄξος, ἀπέπεμψ᾽ ὁ κατάρατος δοὺς ὕδωρ ; Eur.~\textit{Hipp.}~952–953 ἤδη νυν αὔχει καὶ δι᾽ ἀψύχου βορᾶς σίτοις καπήλευ᾽ ; Aesch.~fr.~322~R. (Suid.~κ~337) ὁ δὲ Αἰσχύλος τὰ δόλια πάντα καπηλεῖα καλεῖ· κάπηλα προφέρων τεχνήματα ; Vita Eur.~1 (Σ Eur.~\textit{I,} p.~1 Schwartz) Εὐριπίδης ὁ ποιητὴς υἱὸς ἐγένετο Μνησαρχίδου καπήλου καὶ Κλειτοῦς λαχανοπώλιδος ; Poll.~7.193 (Bethe II, p.~105,~3–5) κάπηλοι δὲ οὐ μόνον οἱ μεταβολεῖς, ἀλλὰ καὶ οἱ τὸν οἶνον κεραννύντες· ὅθεν καὶ Σάραμβον ὁ Πλάτων (Gorg 518B) κάπηλον ὠνόμασεν, ἐπαινῶν αὐτὸν ἐπ᾽ οἰνουργίᾳ ; Poll.~7.194 (Bethe II, p.~105,~10–12) κάπηλιν δ᾽ ἔλεγον, ὡς Ἀριστοφάνης (Plut 436) ἆρ᾽ ἐστὶν ἡ κάπηλις ἡ ἐκ [τῶν] γειτόνων ; Hsch.~π~1552 Πέρδικος καπηλεῖον· χωλὸς κάπηλος ὁ Πέρδιξ ἦν. ἔνθεν ἔνιοι τὴν παροιμίαν φασὶ διαδοθῆναι ; Aesch.~\textit{Sept.}~545 ἐλθὼν δ᾽ ἔοικεν οὐ καπηλεύσειν μάχην ; D.~25.46 εἰ δὲ κάπηλός ἐστι πονηρίας καὶ παλιγκάπηλος καὶ μεταβολεύς, καὶ μόνον οὐ ζυγὰ καὶ στάθμ᾽ ἔχων πάνθ᾽ ὅσα πώποτ᾽ ἔπραξεν ἐπώλει ; Poll.~6.128 (Bethe II, p.~34) Βίοι ἐφ᾽ οἷς ἄν τις ὀνειδισθείη, πορνοβοσκός, κάπηλος ; Synag.~κ~68 (Cunningham, p.~284) κάπηλος· μεταβολεύς, πραγματευτής ; \textit{idem} Suid.~κ~336 (add. κάπηλος ἀσπίδων. ὁ μετάβολος παντὸς πράγματος)}
\lp{347 κάπηλος]}{Ar.~\textit{Lys.}~466 πολλήν γ᾽, ἐάνπερ πλησίον κάπηλος ᾖ ; Ar.~\textit{Thesm.}~737 ὦ μέγα καπήλοις ἀγαθὸν ἡμῖν δ᾽ αὖ κακόν ; Ar.~\textit{Av.}~1292 πέρδιξ μὲν εἷς κάπηλος ὠνομάζετο ; Ar.~\textit{Eccl.}~49 τὴν τοῦ καπήλου δ᾽ οὐχ ὁρᾷς Γευσιστράτην ; Ar.~\textit{Eccl.}~154 ἐν τοῖς καπηλείοισι λάκκους ἐμποιεῖν ; Hsch.~κ~703 καπηλεύει· μεταπωλεῖ, οἰνοπωλεῖ. καὶ τὰ πρὸς τὰς τροφὰς καὶ πόσεις ; Hsch.~κ~705 κάπηλος· μεταπράτης, ὁ τὰ πρὸς τὴν κάπην πιπράσκων· κάπη δὲ ἡ τροφή ; Isoc.~7.49 ἐν καπηλείῳ δὲ φαγεῖν ἢ πιεῖν οὐδεὶς οὐδ᾽ ἂν οἰκέτης ἐπιεικὴς ἐτόλμησεν ; Lys.~1.24 καὶ δᾷδας λαβόντες ἐκ τοῦ ἐγγύτατα καπηλείου εἰσερχόμεθα ; Hsch.~κ~704 καπηλεύοντες· πραγματευόμενοι (2. Ep. Cor. 2, 17) ; Hsch.~κ~707 καπηλοδύτης· ὁ ἐν τῷ καπηλείῳ ἀναστρεφόμενος}
\lp{347 κάπηλος ἢ καπηλίς]}{Ar.~\textit{Plu.}~1120 πρότερον γὰρ εἶχον μὲν παρὰ ταῖς καπηλίσιν ; Hor.~\textit{Sat.}~1.1.29 perfidus hic caupo, miles nautaeque ; Mart.~3.57.1–2 Callidus inposuit nuper mihi copo Ravennae: / Cum peterem mixtum, vendidit ille merum ; Hdt.~1.94.1 πρῶτοι δὲ ἀνθρώπων τῶν ἡμεῖς ἴδμεν νόμισμα χρυσοῦ καὶ ἀργύρου κοψάμενοι ἐχρήσαντο, πρῶτοι δὲ καὶ κάπηλοι ἐγένοντο ; Plu.~\textit{Lys.}~13.5 ἀλλὰ καὶ ὁ κωμικὸς Θεόπομπος ἔοικε ληρεῖν ἀπεικάζων τοὺς Λακεδαιμονίους ταῖς καπηλίσιν, ὅτι τοὺς Ἕλληνας ἥδιστον ποτὸν τῆς ἐλευθερίας γεύσαντες ὄξος ἐνέχεαν ; D.Chr.~31.38 ὅτι μηδὲ τοῖς καπήλοις μηδεὶς ἔτι ῥᾳδίως πρόσεισι, παρ᾽ οἷς ἂν πονηρὰ ᾖ τὰ μέτρα ; Luc.~\textit{Herm.}~59 ὥσπερ οἱ κάπηλοι, κερασάμενοί γε οἱ πολλοὶ καὶ δολώσαντες καὶ κακομετροῦντες ; Lib.~\textit{Or.}~46.10 Ἠδίκουν γάρ, φησίν, οἱ κάπηλοι περὶ τὰ μέτρα τοὺς πίνοντας ; Suid.~κ~334 Καπηλικῶς: ἀντὶ τοῦ πανουργικῶς. ἐπεὶ οἱ κάπηλοι ὀνθυλεύουσι τὸν οἶνον, συμμιγνύντες αὐτῷ σαπρόν ; Suid.~κ~335 Καπηλίς: παρὰ τὸ κακύνειν τὸν πηλόν. πηλὸς δὲ ὁ οἶνος ; Zen.~5.1 Λυδὸς καπηλεύει: φασὶ Κύρον περιγενόμενον τῶν Λυδῶν προςτάξαι αὐτοῖς καπηλεύειν, καὶ μὴ κεκτῆσθαι ὅπλον, ἀλλὰ ποδήρεις ἀμφιέννυσθαι χιτῶνας· καὶ οὕτως εἰς θῆλυ μᾶλλον μεταῤῥυθμισθέντας τοὺς Λυδοὺς ἀγοραίους γενέσθαι καπήλους ; Greg. Cypr.~2.52 Ἔγχος ἐπὶ τῷ ἔγχει: γραῦν φασι τοῦτο εἰπεῖν καπήλῳ ἐνεχυράζουσαν ἔγχος ; Ar.~\textit{Plu.}~1063 οὐ δῆτ᾽, ἐπεὶ νῦν μὲν καπηλικῶς ἔχει ; Theopomp. Com.~\textit{Καπηλίδες} (Kock I, p.~739) Θεόπομπος Καπήλισι ; Hor.~\textit{Sat.}~1.5.4 differtum nautis cauponibus atque malignis ; Mart.~1.56.1–2 Continuis vexata madet vindemia nimbis: / Non potes, ut cupias, vendere, copo, merum ; Is.~1.22 LXX (ap. Cyr. Alex. In Am.~1.34) οἱ κάπηλοί σου μίσγουσι τὸν οἶνον ὕδατι ; Hdt.~1.155.4 πρόειπε δ᾽ αὐτοῖσι κιθαρίζειν τε καὶ ψάλλειν καὶ καπηλεύειν παιδεύειν τοὺς παῖδας ; Pl.~\textit{Lg.}~918a κιβδήλοις δ᾽ ἐπιτηδεύμασιν ἕπεται κατὰ πόδα καπηλείας ἐπιτηδεύματα ; Zen.~4.44 Μέμνηται ταύτης Θεόπομπος ἐν Καπηλίσιν ; Erasm.~\textit{Adag.}~2.6.96 Lydus cauponatur. Λυδὸς καπηλεύει. i. e. Lydus cauponatur. In effoeminatos \& voluptatibus addictos jaciebatur … Herodotus lib. 1. tradit, Lydos primos omnium caupones \& institutores exstitisse ; Ar.~\textit{Thesm.}~630–631 φέρ᾽ ἴδω, τί μέντοι πρῶτον ἦν; ἐπίνομεν. / τί δὲ μετὰ τοῦτο δεύτερον; προὐπίνομεν ; Pl.~\textit{Cra.}~417b οὐχὶ καθάπερ οἱ κάπηλοι αὐτῷ χρῶνται ; Ath.~14.656f ὥσπερ αἱ καπηλίδες τὰ δελφάκια τρέφουσιν}
\lp{347 κάπηλος … τοῦ χοός]}{Eub.~fr.~80~Kock (Ath.~11.473e) ἐγὼ δὲ καὶ γὰρ ἔτυχεν ὂν καταντικρὺ / τῆς οἰκίας καινὸν καπηλεῖον μέγα / ἐνταῦθ᾽ ἐπετήρουν τὴν τροφὸν τῆς παρθένου, / κεράσαι κελεύσας τὸν κάπηλόν μοι χόα / ὀβολοῦ ; EM~490.4–14 Κάπηλος· ὁ μετάβολος καὶ οἰνοπώλης· παρὰ τὸ χέειν τὸν πηλόν, ἤγουν τὸν οἶνον … πόθεν δὲ ἐν τῷ καπήλῳ ὁ οἶνος ἔγχειται; καὶ γὰρ ἕτερόν ἐστιν οἰνοπώλης, καὶ ἕτερον καπηλεῖον. καθόλου γὰρ τοὺς πωλοῦντάς τι καπήλους ἔλεγον}
\lp{347 τοῦ χοός]}{Ar.~\textit{Thesm.}~746 πόσ᾽ ἔτη δὲ γέγονε; τρεῖς Χοᾶς ἢ τέτταρας ; Ar.~\textit{Eccl.}~44 τὴν ὑστάτην ἥκουσαν οἴνου τρεῖς χοᾶς ; Σ Ar.~\textit{Thesm.}~347 (Dübner, p.~266,~51–52) τοῦ διξέστου. χοῦς γὰρ δύο ξέσται, χοεὺς δὲ ἕξ ; Ar.~\textit{Eq.}~95 ἀλλ᾽ ἐξένεγκέ μοι ταχέως οἴνου χοᾶ ; Ar.~\textit{Nub.}~1238 ἓξ χοᾶς χωρήσεται ; Ar.~\textit{Pax}~537 δούλης μεθυούσης, ἀνατετραμμένου χοῶς ; Hsch.~χ~573 χόα· χοῦν. μέτρον ὑγρῶν ; Hdn.~\textit{Gr.}~2.706.1–7 Lentz (Choerob.~240.33) τὸ χόα, ἀπὸ τοῦ χοῦς χοός … τὸ δὲ παρὰ τῷ Μενάνδρῳ χοᾶ … ἀπὸ τοῦ χοεύς χοέος … γέγονε χοᾶ κατὰ κρᾶσιν ; Diogenian.~2.83b Ἀπὸ χοέως σπάσον: παρόσον οἱ μεθύοντες ἀληθεύουσιν ; \textit{idem} Apostol.~3.67 (Ἀπ᾽ οἰνοχοέως) ; App. Prov.~5.29 Χοᾶς ἓξ χωρήσει: τουτέστι ὁ ἐγκέφαλος αὐτοῦ ; Cratin.~fr.~199.3–4~K.-A. (187~Kock) ἐγᾦδα· συντρίψω γὰρ αὐτοῦ τοὺς χοᾶς / καὶ τοὺς καδίσκους συγκεραυνώσω σποδῶν}
\lp{347–348 τοῦ χοὸς ἢ τῶν κοτυλῶν]}{Poll.~7.194–195 (Bethe II, p.~105,~14–21) καὶ τὸ κοτυλίζειν εἴρηται μὲν καὶ ἐπὶ τῶν ἄλλων τῶν κατὰ μέρος πιπρασκομένων ἀλλ᾽ οὐκ ἀθρόως … ἴδιον δ᾽ ἂν εἴη τῶν τὸν οἶνον πιπρασκόντων. καὶ Φερεκράτης μὲν εἴρηκε «μηδὲν κοτυλίζειν, ἀλλὰ καταπάττειν χύδην», Ἀριστοφάνης δὲ «κιρνάντες γὰρ τὴν πόλιν ἡμῶν κοτυλίζετε τοῖσι πένησιν» ; Ar.~\textit{Plu.}~737 καὶ πρίν σε κοτύλας ἐκπιεῖν οἴνου δέκα}
\lp{347–348 τοῦ χοὸς ἢ τῶν κοτυλῶν τὸ νόμισμα]}{Arist.~\textit{Ath.}~51.2 κληροῦνται δὲ καὶ μετρονόμοι ⟨ι⟩, πέντε μὲν εἰς ἄστυ, ε δὲ εἰς Πειραιέα. καὶ οὗτοι τῶν μέτρων καὶ τῶν σταθμῶν ἐπιμελοῦνται πάντων, ὅπως οἱ πωλοῦντες χρήσονται δικαίοις ; Thphr.~\textit{Char.}~30.11 καὶ τὰ τοιαῦτα· Φειδωνείῳ μέτρῳ τὸν πύνδακα εἰσκεκρουμένῳ μετρεῖν αὐτὸς τοῖς ἔνδον τὰ ἐπιτήδεια σφόδρα ἀποψῶν ; Diogenian.~5.69 Κωδάλου χοῖνιξ: ἐπὶ τῶν μεγάλοις μέτροις κεχρημένων ; Arist.~\textit{Ath.}~10.2 ἐγένετο καὶ τὰ μέτρα μείζω τῶν Φειδωνείων}
\lp{348 τὸ νόμισμα διαλυμαίνεται]}{D.~20.167 θαυμάζω δ᾽ ἔγωγε, εἰ τοῖς μὲν τὸ νόμισμα διαφθείρουσιν θάνατος παρ᾽ ὑμῖν ἐστιν ἡ ζημία ; D.~24.212 νόμος ἐστὶν ἁπάσαις, ὡς ἔπος εἰπεῖν, ταῖς πόλεσιν, ἐάν τις τὸ νόμισμα διαφθείρῃ, θάνατον τὴν ζημίαν εἶναι ; Pl.~\textit{Lg.}~916d ὁ δὲ ἀλλαττόμενος ἢ νόμισμα ἀντὶ νομίσματος … ἀκίβδηλον πᾶν διδότω καὶ δεχέσθω τῷ νόμῳ συνεπόμενος … κιβδηλείαν δὲ χρὴ πάντα ἄνδρα διανοηθῆναι καὶ ψεῦδος καὶ ἀπάτην ὡς ἕν τι γένος ὄν ; Hsch.~π~664 παράσημον· ἀδόκιμον, κίβδηλον (Aesch. Ag. 781. Ar. Ach. 518. cf. Eur. Hipp. 1115). ἀπὸ τῶν [καὶ] νομισμάτων ⟨ὧν⟩ τὰ κίβδηλα παράσημα λέγεται (Dem. 24, 213. cf. 18, 242) ; Hsch.~α~2396 ἀκίβδηλον· καθαρόν, ἁγνόν, σεπτόν. ἢ δόκιμον, ἀπαρεγχείρητον, ἐπὶ νομίσματος}
\lp{348 τὸ νόμισμα]}{Ar.~\textit{Nub.}~247–248 πρῶτον γὰρ θεοὶ ἡμῖν νόμισμ᾽ οὐκ ἔστι ; Ar.~\textit{Av.}~1040–1041 χρῆσθαι Νεφελοκοκκυγιᾶς τοῖσδε τοῖς μέτροισι καὶ σταθμοῖσι καὶ ψηφίσμασι ; Ar.~\textit{Ran.}~720–721 ἔς τε τἀρχαῖον νόμισμα καὶ τὸ καινὸν χρυσίον. οὔτε γὰρ τούτοισιν οὖσιν οὐ κεκιβδηλευμένοις ; Suid.~ν~466 Νόμισμα: δύο σημαίνει· τὸ νόμιμον ἔθος· ὡς Ἀριστοφάνης ἐν Νεφέλαις, ἡμῖν νόμισμ᾽ οὐκ ἔστι· καὶ τὸ κόμμα τοῦ τετυπωμένου χαλκοῦ ; Eur.~fr.~542~N.² οὔτοι νόμισμα λευκὸς ἄργυρος μόνον / καὶ χρυσός ἐστιν, ἀλλὰ κἀρετὴ βροτοῖς / νόμισμα κεῖται πᾶσιν ; Aesch.~\textit{Sept.}~269 Ἑλληνικὸν νόμισμα θυστάδος βοῆς ; Soph.~\textit{Ant.}~295–296 οὐδὲν γὰρ ἀνθρώποισιν οἷον ἄργυρος / κακὸν νόμισμ᾽ ἔβλαστε ; Arist.~\textit{EN}~1164a1 ἐνταῦθα μὲν οὖν πεπόρισται κοινὸν μέτρον τὸ νόμισμα ; Harp. s.v.~διάμετρον ΔΙΑΜΕΤΡΟΝ τὸ ἐλλεῖπον ἀπὸ τοῦ δικαίως μετρηθέντος ; App. Prov.~1.81 Γνῶθι σαυτόν καὶ Τὸ νόμισμα παραχάραξον: παραγγέλματα Πυθικά ; Antiph.~fr.~40~Kock ὅστις δὲ μεῖζον ἢ κατ᾽ ἄνθρωπον φρονεῖ / . . . μικρῷ πεποιθὼς ἀθλίῳ νομίσματι ; Eur.~\textit{Cyc.}~160 πρὸς τῷδε μέντοι καὶ νόμισμα δώσομεν}
\lp{348 διαλυμαίνεται]}{Ar.~fr. inc. ap. Harp. s.v.~Κεῖοι Ἀριστοφάνης δὲ περὶ Σικελίας φησὶν οἵα πόλις τάλαινα διαλυμαίνεται ; Eur.~\textit{Hipp.}~1349 χρησμοῖς ἀδίκοις διελυμάνθην ; Ar.~\textit{Ran.}~59 τοιοῦτος ἵμερός με διαλυμαίνεται ; Ar.~\textit{Ran.}~1062 ἁμοῦ χρηστῶς καταδείξαντος διελυμήνω σύ ; Eur.~\textit{Or.}~1515 οὐ γάρ, ἥτις Ἑλλάδ᾽ αὐτοῖς Φρυξὶ διελυμήνατο ; Isoc.~4.110 καὶ τὰς αὑτῶν πατρίδας διαλυμηνάμενοι ; Hsch.~λ~1407 λυμαίνεσθαι· καθαίρειν. φθείρειν ; Hsch.~λ~1408 λυμαίνεται· διαφθείρεται (Eur. Bacch. 354) ; Hsch.~λ~1482 λωβᾶται· ἐξαπατᾷ. λυμαίνεται ; Hdt.~9.112.1 ἡ Ἄμηστρις μεταπεμψαμένη τοὺς δορυφόρους τοῦ Ξέρξεω διαλυμαίνεται τὴν γυναῖκα τοῦ Μασίστεω ; D.~9.31 τὰ μὴ προσήκοντ᾽ ἀπώλλυε καὶ ἐλυμαίνετο ; D.~\textit{Ep.}~1.12 τὰ γὰρ ὀρθῶς βουλευθέντα … τῷ τοὺς ἐπιστάντας ἄλλως χρήσασθαι διελυμάνθη ; Thphr.~\textit{CP}~2.1.6 ἐπιγινόμενα ψύχη διελυμήνατο καὶ ἀπέκαυσεν ; Suid.~λ~838 Λυμαίνομαί σοι· δοτικῇ· αἰτιατικῇ δέ· ἐλυμήνατο αὐτὴν ὗς ἐκ δρυμοῦ}
\lp{348 τῶν κοτυλῶν]}{Hsch.~χ~657 χοῦς· ἄνθρωπος (Gen. 2, 7). καὶ τὸ ἐπιβαλλόμενον τῷ ὀρόφῳ χῶμα (Thuc. 2, 76, 2 . .). καὶ μέτρον ⟨ηʹ⟩ κοτυλῶν (Ar. Equ. 95 . .). καὶ ἐρ⟨ε⟩ίπιον ; Hsch.~κ~3813 κοτύλη· εἶδος μέτρου. καὶ τρώγλη. ἢ τὸ κοῖλον τῆς χειρός, καὶ πᾶν κοῖλον ; Suid.~κ~2166 Κοτύλη: ὃ νῦν καλοῦσιν ἡμίξεστον … εἶδος ποτηρίου ὁμωνύμου τῷ ἀγγείῳ τῷ ὑγρῷ, ὡς χοῖνιξ. καὶ εἶδος μέτρου ; Hsch.~κ~3812 κοτύλαια· μικρά· παρὰ τὴν κοτύλην ; Synag.~κ~439 (Cunningham, p.~308) κοτύλη· εἶδος μέτρου. καὶ τὸ κοῖλον τοῦ ὀστέου, ἔνθα ἡ κεφαλὴ τοῦ μηροῦ ἐνστρέφεται}
\lp{348–351 διαλυμαίνεται … κακῶς ἀπολέσθαι … πολλὰ δοῦναι κἀγαθά]}{SIG³~997.1–16 [ἰ]χθῦς ἱεροὺς μὴ ἀδικεῖν μηδὲ σκεῦος τῶν τῆς θεοῦ λυμαίνεσθαι … ὁ τούτων τι ποιῶν κακὸς κακῇ ἐξωλείᾳ ἀπόλοιτο, ἰχθυόβρωτος γενόμενος … τοῖς δὲ συμφυλάσσουσιν καὶ ἐπαύξουσιν τὰ τῆς θεοῦ τίμια … βίου καὶ ἐργασίας καλῆς γένοιτο παρὰ τῆς θεοῦ ὄνησις}
